%!TEX program = lualatex
\documentclass[10pt, openany]{book}

\usepackage{../preamble}
\addbibresource{../ref.bib}

\fancyhead[LE,RO]{\thepage}
\fancyhead[RE]{\textbf{几何讲义}}
\fancyhead[LO]{\textit{范畴论观点}}

\begin{document}

\frontmatter
\title{\textbf{几何讲义}\\  }
\author{Lao Chen}
\date{\today}
\maketitle
\tableofcontents

\chapter*{符号列表}
\addcontentsline{toc}{chapter}{符号列表}
\glsaddall 
\printglossaries

\newglossaryentry{cat-set}{
  name={\ensuremath{\textbf{Set}}},
  description={集合范畴}}

\newglossaryentry{cat-top}{
  name={\ensuremath{\textbf{Top}}},
  description={拓扑空间范畴}}

\newglossaryentry{cat-cring}{
  name={\ensuremath{\textbf{CRing}}},
  description={交换环范畴}}

\newglossaryentry{contra-variant-hom}{
  name={\ensuremath{h^A = \text{Hom}(-,A)}},
  description={反变Hom函子}}

通常：

大写字母$X,Y,Z$表示集合、拓扑空间、流形；
$A,B,C$为交换环或交换代数．
小写字母为其中的元素．

$\mathbb Z,\mathbb Q,\mathbb R,\mathbb C$表示数环、数域，
$k$表示交换域，
$\mathbb Z[\sqrt{n}]$为二次扩环．

$n\in\mathbb Z$表示整数，
$p$表示非零素数．
$(n)$表示理想，
$\mathfrak{p}$表示素理想，
$\mathfrak{m}$表示极大理想．

\mainmatter

\chapter{点集拓扑}

\section{从集合构造幂集}
\begin{framed}
\noindent\textbf{定义（幂集）：}给定集合 $X$，其所有的子集构成的集合称为\textbf{幂集}（Power Set），记作 $\mathcal{P}(X)$ 或 $2^X$．
\end{framed}
在拓扑学中，空间的拓扑结构 $\tau$ 本质上就是幂集 $\mathcal{P}(X)$ 的一个特定子集．

\section{拓扑公理}
\begin{framed}
\noindent\textbf{定义（拓扑空间与开集）：}一个集合 $X$ 上的拓扑 $\tau \subseteq \mathcal{P}(X)$ 是一族称为“\textbf{开集}”的集合，必须满足以下三个公理：
\begin{enumerate}
    \item \textbf{平凡性}：空集 $\emptyset$ 和全集 $X$ 都属于 $\tau$．
    \item \textbf{有限交封闭}：任意有限个开集的交集仍属于 $\tau$．
    \item \textbf{任意并封闭}：任意多个开集的并集仍属于 $\tau$．
\end{enumerate}
满足这些公理的二元组 $(X, \tau)$ 称为\textbf{拓扑空间}．
\end{framed}

\section{平凡拓扑、离散拓扑与拓扑结构的偏序包含}
\begin{framed}
\noindent\textbf{定义（平凡拓扑与离散拓扑）：}在同一个集合 $X$ 上可以赋予极其不同的拓扑结构．包含开集最少的拓扑 $\tau = \{\emptyset, X\}$ 称为\textbf{平凡拓扑}（Trivial Topology）；所有子集都是开集的拓扑 $\tau = \mathcal{P}(X)$ 称为\textbf{离散拓扑}（Discrete Topology）．
\end{framed}
如果集合 $X$ 上有两个拓扑 $\tau_1$ 和 $\tau_2$，且 $\tau_1 \subseteq \tau_2$，我们称 $\tau_2$ 比 $\tau_1$ \textbf{更细}（finer），或称 $\tau_1$ 比 $\tau_2$ \textbf{更粗}（coarser）．所有可能的拓扑结构通过包含关系构成了一个偏序集．

\section{无限并，从区间套极限得到单点，单点是闭集的意义}
\begin{framed}
\noindent\textbf{定义（闭集）：}在拓扑空间中，作为开集补集的子集称为\textbf{闭集}．
\end{framed}
拓扑公理保证了开集对“无限并”是封闭的，但对“无限交”不一定封闭．相对地，闭集对任意交都是封闭的．利用闭集的无限交特性，我们可以通过“区间套极限”的方式将空间逼近到一个单点．在 $T_1$ 空间中，每一个单点集 $\{x\}$ 都是闭集．如果一个拓扑空间中所有的单点集不仅是闭集，同时也是开集，那么这个空间的拓扑结构必定退化为\textbf{离散拓扑}．

\section{Zariski拓扑：用零点集构造}
\begin{framed}
\noindent\textbf{定义（Zariski 拓扑）：}设 $k$ 为域，对于多项式环 $k[x_1, \dots, x_n]$ 中的子集 $S$，定义其公共零点集为：
$$ V(S) = \{ \mathbf{a} \in k^n \mid f(\mathbf{a}) = 0, \forall f \in S \} $$
规定所有的 $V(S)$ 为空间中的\textbf{闭集}，其补集即为开集．这种由此定义的拓扑结构称为 \textbf{Zariski 拓扑}．
\end{framed}
在代数几何中，空间形体的构造由代数方程主导．Zariski 拓扑是一种极其“粗糙”的拓扑，在不可约代数簇上，任意两个非空开集必定相交．

\section{Hausdorff空间}
\begin{framed}
\noindent\textbf{定义（Hausdorff 空间）：}拓扑空间 $(X, \tau)$ 称为 \textbf{Hausdorff 空间}（或 $T_2$ 空间），如果对于空间中任意两个不同的点 $x, y \in X, x \neq y$，都存在两个不相交的开邻域 $U$ 和 $V$，使得 $x \in U, y \in V$ 且 $U \cap V = \emptyset$．
\end{framed}
这一分离性质是保证极限唯一的基石，也是分析学得以展开的前提．
\chapter{交换代数}

定理：
\textbf{理想是素的 iff. 商环是整环．
理想是极大的 iff. 商环是域．}

$$
% https://tikzcd.yichuanshen.de/#N4Igdg9gJgpgziAXAbVABwnAlgFyxMJZAJgBoAGAXVJADcBDAGwFcYkQBBAHS6zB5wwAHjgBGAM2ABhAEp8A5gF8Qi0uky58hFGQCM1Ok1bseAW3o4AFuIBO9ANYACNI559XXQSOABlNDABjRQ9HDhU1EAxsPAIiABZSfRoGFjZETgB6MwtrOycXNzAPLzFJAEkAEWVVdWitInJEgxTjdI4AfWyrWwdnRwBeRwAKbi44GBxTPmY4D3NuvOcASgA9YABaXUUwmsiNGO0SUmJmozSQeyGu3N60JYHQzq55m-ys55ye-Mfrr76ebCmGAAR2KwhwwAAYnYgiEOO8Xn80CoDDAoPJ4ERQLYIKYkABmGg4CBIBIgHD0LCMdgUqnhbE2XFIRrkkmIQnkynU9KWCAQez0kA4vGIFnEpBkTl09K06m7YUSolssmy9i8-koxRAA
\begin{tikzcd}
                                                                 &  & A\in\textbf{CRing} \arrow[lld, hook] \arrow[rrd, tail]                                      &  &                                                 \\
A_\mathfrak p = (A\setminus \mathfrak p)^{-1}A \arrow[rrd, tail] &  & \mathfrak p \in \text{Spec} \ A                                                             &  & A/\mathfrak p \in \textbf{ID} \arrow[lld, hook] \\
                                                                 &  & k(\mathfrak p) = A_\mathfrak p/\mathfrak p A_\mathfrak p \simeq \text{Frac} \ A/\mathfrak p &  &                                                
\end{tikzcd}
$$


$$
% https://tikzcd.yichuanshen.de/#N4Igdg9gJgpgziAXAbVABwnAlgFyxMJZAJgBoAGAXVJADcBDAGwFcYkQBBAHS6zB5wwAHjgBGAM2ABhAEp8A5gF8Qi0uky58hFGQCM1Ok1bseAW3o4AFuIBO9ANYACU4559XXQSOABZekMUPRw4VNRAMbDwCIgAWUn0aBhY2RE4AejMLazsnFzcwDy8xSQAxRihlVXVIrSJyeIMk41SOAH1Mq1sHZ0cAXkcACm4uOBgcUz5mOA9zTpznAEoAPWAAWl1FEKrwjSjtElJiRqMUznauWezu0wyLrK7c4PPLh56ebFMYAEdg25f50wqAwwKDyeBEUC2CCAxAAZhoOAgSDiIBw9CwjHYaIxoUhNmhSHqqKRcIR6MxqUsEAg9lxIChMKJiKQZFR5Kx7O2DJZCJJKOxFJAVJpQMUQA
\begin{tikzcd}
                                                                 &  & A\in\textbf{CRing} \arrow[lld, hook] \arrow[rrd, tail]       &  &                                                  \\
A_\mathfrak m = (A\setminus \mathfrak m)^{-1}A \arrow[rrd, tail] &  & \mathfrak m \in \text{Max} \ A                               &  & A/\mathfrak m \in \textbf{Fld} \arrow[lld, hook] \\
                                                                 &  & A_\mathfrak m/\mathfrak m A_\mathfrak m \simeq A/\mathfrak m &  &                                                 
\end{tikzcd}
$$

$A/\mathfrak p\in\textbf{ID}$是整环，
可做分式域$\text{Frac} \ A/\mathfrak p$．
$A\smallsetminus \mathfrak p$是乘法系统，
从而有局部化$A_\mathfrak p$，
及其不可逆元的全体构成的唯一的极大理想$\mathfrak pA_\mathfrak p$，
进而有剩余域$k(\mathfrak p) = A_\mathfrak p/\mathfrak pA\mathfrak p$．
有典范同构
$k(\mathfrak p) = A_\mathfrak p/\mathfrak pA\mathfrak p
\simeq
\text{Frac} \ A/\mathfrak p$

对于极大理想$\mathfrak m\in\text{Max} \ A$，
$A/\mathfrak \in \textbf{Fld}$已经是域，
$\text{Frac} \ A/\mathfrak m \simeq A/\mathfrak m$:
有典范同构
$k(\mathfrak m) = A_\mathfrak m/\mathfrak m A_\mathfrak m \simeq \text{Frac} \ A/\mathfrak m \simeq A/\mathfrak m$

整数环：
以$(3) = \text{Max} \ \mathbb Z$为例，
局部环$\mathbb Z_{(3)} = \{a/s \mid a,s\in\mathbb Z,\ s \notin (3)\}$
有唯一极大理想$(3)\mathbb Z_{(3)} = \{3a/s \mid a,s\in\mathbb Z,\ s \notin (3)\}$，
包含所有不可逆元．
剩余域$\mathbb Z_{(3)}/(3)\mathbb Z_{(3)} \simeq \mathbb Z/(3)$，
正是Galois域$\mathbb F_3$．

多项式环：
以$((x^2 + 1)) = \text{Max} \ \mathbb R[x]$为例，
局部环$\mathbb R[x]_{((x^2 + 1))} = \{p/q \mid p,q\in\mathbb R[x],\ q \notin ((x^2 + 1))\}$
有唯一极大理想$((x^2 + 1))\mathbb R[x]_{((x^2 + 1))} = \{(x^2 + 1)a/s \mid p,q\in\mathbb R[x],\ s \notin ((x^2 + 1))\}$，
包含所有不可逆元．
剩余域$\mathbb R[x]_{((x^2 + 1))}/((x^2 + 1))\mathbb R[x]_{((x^2 + 1))} \simeq \mathbb R[x]/((x^2 + 1))$．
对于域$k$和不可约多项式$p \in k[x]$，
商环$k[x]/(p)$是$k$的扩域．
在这里正是复域$R[x]/((x^2 + 1)) \simeq \mathbb C$．

函数环：
以$\mathfrak m_x = \{f\in C(X) \mid f(x) = 0\} \in \text{Max} \ C(X)$为例，
这是在$x \in X$点退化的函数全体构成的极大理想．
局部环$C(X)_{\mathfrak m_x} = \{f/g \mid f,g\in C(X),\ g(x) \neq 0 \}$，
若$X$是仿紧Hausdorff空间则它同构于函数芽环$\mathcal O_x$．
有唯一极大理想$\mathfrak m_x C(X)_{\mathfrak m_x} = \{f/g \mid f,g\in C(X),\ f(x) = 0,\ g(x) \neq 0 \}$，
包含所有不可逆元，
即所有在$x$退化的函数芽．
剩余域$C(X)_{\mathfrak m_x}/\mathfrak m_x C(X)_{\mathfrak m_x} \simeq \mathbb R$．

若$A\in\textbf{CC}^*$还是交换C*-代数，
其所有素理想都是极大理想．
作为集合
$\text{Max} \ A = \text{Spec} \ A$．
此时给$\text{Max} \ A$赋予的是Gelfand拓扑，
从而有函子$\text{Max}: \textbf{CC}^* \to \textbf{Set}$．

\section{交换环的理想}


\begin{framed}
\noindent
$\forall a,b \in A \in \textbf{CRing}$，
$(ab=0) \Rightarrow (a=0) \vee (b=0)$，
则称 $A$ 为\textbf{整环}（Integral Domain）．
\end{framed}

整数环 $\mathbb{Z}$ 是整环；而商环 $\mathbb{Z}/(n)$ 不是．
商环的元素记为$[a]_n = a + (n) \in \mathbb{Z}/(n)$，
例如$n=6$时，
$[2]_6 \times [3]_6 = [2 \times 3]_6 = 0 \in \mathbb{Z}/(6)$，
它包含零因子．

\begin{framed}
\noindent\textbf{定义（素理想与极大理想）：}
\textbf{素理想}（Prime Ideal） $\mathfrak{p}$ 要求：若 $ab \in \mathfrak{p}$，则必定有 $a \in \mathfrak{p}$ 或 $b \in \mathfrak{p}$．
\textbf{极大理想}（Maximal Ideal） $\mathfrak{m}$ 则是真包含于 $A \in \textbf{CRing}$ 且不可被其他真理想真包含的理想．
\end{framed}

$A$的素理想的全体构成素谱$\text{Spec} \ A$，
后面我们将了解它具有几何意义，
可以视为拓扑空间$X = \text{Spec} \ A$，
而素理想$\mathfrak p\in \text{Spec} \ A$视为这个空间的点．

$A$的极大理想的全体构成Gelfand谱$\text{Max} \ A$，
当$A$是含幺交换C*-代数时，
它在Gelfand拓扑（弱*拓扑）下构成紧Hausdorff空间，
这是算子代数重要的基础．

\section{扩环}

以有限集合$X = \{x_i\} \in\textbf{FinSet}$为基，
以集合$A$为系数集，
有函数集$h^A X = A^X = \textbf{Set}(X,A)$，
给定其中的函数$f:X \to A$为每个$x_i$指定了系数$a_i = f(x_i)$，
以及对应的$(a_i,x_i)$，
从而产生形式和：
$\sum_i a_i x_i = a_1 x_1 + a_2 + x_2 + \cdots = \simeq \{ (a_1,x_1),(a_2,x_2), \cdots \}_i$，
这里$a_1 x_1 + a_2 + x_2 + \cdots$只是形式加法，
并不意味着有可加性．
形式和的全体记为：
$A[X] = \{ \sum_i a_i x_i \} \simeq h^A X$．

若生成集$X = \textbf 1 = \{1\}$为单点集，
$\sum_i a_i x_i = a_1 1$，
因此
$A[\textbf{1}] = \{ a_1 1 \} \simeq \{ a_1 \} \simeq A \simeq h^A \textbf 1$，
形式和的全体相当于系数集．
$\mathbb Z$可以视为基$X=\{1\}$的整系数形式和的集合．

当系数集$A \in \textbf{Ab}$是Abel群，
借助$A$中的交换加法，
定义形式和$A[X]$的含幺可逆交换加法：
$\sum_i a_i x_i + \sum_i b_i x_i = \sum_i (a_i + b_i) x_i$，
使得$A[X] \in \textbf{Ab}$构成Abel群．

有限集合$X$生成的是有限形式和，
若系数$A \in \textbf{CRing}$是交换环，
借助$A$中的交换乘法，
定义形式和$A[X]$的含幺交换乘法：
$(\sum_i a_i x_i) \cdot (\sum_i b_i x_i)$，
展开合并后就会出现生成元的有限次幂．
为了让$A[X]$的乘法封闭，
需要对幂进行处理．
以下固定整数环系数$A = \mathbb Z$，
且明确生成元$1 \in \mathbb Z$是整数，
于是整数$1$作为生成元的任意幂还是它本身，
这样解决了形式和$\mathbb Z[\{1\}]$的乘法封闭性问题，
使得$\mathbb Z[\{1\}] \in\textbf{CRing}$构成交换环，
它就是整数环．

添加一个生成元$s$,
对$X = \{1,s\}$做整系数形式和，
$\sum_i a_i x_i = a_1 1 + a_2 s \simeq \{ (a_1,1),(a_2,s) \}_i$，
形式和的全体为：
$\mathbb Z[X] = \{ a_1 1 + a_2 s \}$．
根据以上讨论让生成元$1$为整数，
为了让形式和$\mathbb Z[X]$的乘法封闭，
要求包含$s$的幂可以回到形式和$a_1 1 + a_2 s$中，
简单方法是要求$s\cdot s = d \mathbb Z$，
这样$(a_1 s) \cdot (a_2 s) = (a_1 a_2)d$，
实现了形式和乘法的封闭性．
记$s = \sqrt d$，
形式和相应简记为：
$a + b\sqrt{d} = a 1 + b \sqrt{d}$．

注意这里还需要考虑根号运算本身的问题：
$d = 24$具有平方因数$2^2$，
导致$\sqrt{2^2 \times 6} = 2 \sqrt 6$，
这样会引入新的生成元$\sqrt 6$，
破坏了封闭性．
我们要求$d$是一个无平方因数，
即$d$的因数分解后没有平方以上的项．
因此我们应当使用无平方因数的根式作为生成元$\sqrt d$．
通过给$\mathbb Z$添加无平方因数的根$\sqrt d$，
得到的环称为\textbf{二次扩环}（Quadratic Extension）
$\mathbb{Z}[\sqrt{d}] = \{a + b\sqrt d\}$．
当 $d=-1$ 时，
形式化的记$ \text i = \sqrt{-1}$，
得到虚二次扩环 $\mathbb{Z}[i] = \{a + b\text i\}$，
即\textbf{Gauss整数环}．


\section{分式}

\begin{framed}
\noindent
交换环$A = (A,+,0,\cdot,1) \in \textbf{CRing}$中，
包含乘法单位元且对乘法封闭的子集称为\textbf{乘法系统}（Multiplicative System）．
\end{framed}

乘法系统$S$的作用是构造分式中的分母．
我们希望分式作为等价类，
例如$2/6 \simeq 1/3$，
基于$2 \times 3 = 1 \times 6$．

\begin{framed}
\noindent
在$A \times S$上构造等价关系
$(a_1,s_1) \sim (a_2,s_2) \Leftrightarrow \exists t \in S:\ t(a_1s_2 - a_2s_1) = 0$，
等价类记为$a/s=[(a_1,s_1)]=[(a_2,s_2)]$，
商环$S^{-1}A = (R \times S)/\sim$称为$A$对$S$的局部化分式环．
\end{framed}

有自然同态$A \to S^{-1}A :: a \mapsto a/1$．
定义中的$t$因子是为了在非整环的情况下，
让等价关系保持传递性，
允许我们消去非整环中零因子带来的歧义．
在整环的情况下可以取$t=1$，
从而退化为简单的$a_1s_2 = a_2s_1$．

局部化给环添加了分母，
使作为分母的乘法系统$S$的元在局部化后称为可逆元．
以下是一些常见的乘法系统及其局部化：

对于 $f \in A$，其所有非负次幂构成的集合 $S = \{1, f, f^2, \dots\}$ 是乘法系统．
例如 $\mathbb{Z}$ 中以 $2$ 的幂次构造的系统．
进而有局部环$A_f = S^{-1}A$．

取$A$的所有非零元$S = A \setminus \{0\}$，
若$A$是整环则$S$是乘法系统：

\begin{framed}
\noindent
整环 $A$ 对 $S = A \setminus \{0\}$ 进行局部化得到的环$\text{Frac} \ A = (A \setminus \{0\})^{-1} A $，
称为\textbf{分式域}（Field of Fractions）．
\end{framed}

$\text{Frac} \ A$它是包含 $A$ 的最小的域．

\begin{itemize}
    \item 整数环 $\mathbb{Z}$ 的分式域是有理数域 $\mathbb{Q}$．
    \item 代数整数环（如 $\mathbb{Z}[i]$）的分式域是对应的代数数域（如 $\mathbb{Q}(i)$）．
    \item 多项式环 $k[x_1, \dots, x_n]$ 的分式域是全体有理函数构成的有理函数域 $k(x_1, \dots, x_n)$．
\end{itemize}


\begin{framed}
\noindent
若一个交换环包含唯一的极大理想，称其为\textbf{局部环}（Local Ring）．
\end{framed}

若 $\mathfrak p \in \text{Spec} \ A$ 是素理想，
则 $S = A \setminus \mathfrak p$ 是乘法系统，
$A_\mathfrak p = S^{-1}A = \{ a/s \mid a \in \mathbb Z,\ s \notin \mathfrak p \}$是局部环．
它有唯一的极大理想$\mathfrak p A_\mathfrak p$．


对于素数$p$，
记$\mathfrak p = (p)$，
$\mathbb{Z}_\mathfrak p$ 是局部环，
它唯一的极大理想是$\mathfrak p \mathbb Z_\mathfrak p$．
例如素数 $3$ 使得所有不能被 $3$ 整除的整数构成乘法系统
$S = \mathbb{Z} \setminus (3)$，
以及局部环 $\mathbb{Z}_{(3)} = \{ a/s \mid a \in \mathbb Z,\ 3 \nmid s \}$．
在这个环中，所有不被 $3$ 整除的数都成为了可逆元．
若$a/s \in \mathbb{Z}_{(3)}$可逆，
则逆元$s/a$需要使得$3 \nmid a$．
那么对于任何$r\in (3)$，
$r/s$的逆元$s/r$无法满足$3 \nmid r$．
于是 $(3)\mathbb{Z}_{(3)} = \{ r/s  \mid r \in (3),\ 3 \nmid s \}$ 作为不可逆元构成的集合，
成为了该局部环的唯一极大理想．


\section{素谱}

现在把前面讨论过的素谱$\text{Spec} \ A$从集合提升为拓扑空间．
令$I \subseteq A$为理想，
素理想$\mathfrak p \in\text{Spec} \ A$，
考虑当它被视为$A$的子集，
包含理想$\mathfrak p \supset I$的情况．

\begin{framed}
\noindent
理想 $I$ 的退化集为包含 $I$ 的所有素理想
$ V(I) = \{ \mathfrak{p} \in \operatorname{Spec}(R) \mid I \subseteq \mathfrak{p} \} $．
\end{framed}

在整数环中，
$I = (60)$，
$\mathfrak p_1 = (2),\ \mathfrak p_2 = (3),\ \mathfrak p_3 = (5)$
都满足这种包含关系．
这里$V(I) = V((60)) = \{ \mathfrak p_1,\ \mathfrak p_2,\ \mathfrak p_3 \}$，
它包含所有整除 $60$ 的素数生成的理想．

退化集$V(I) \subseteq \text{Spec} \ A$以素理想为点，
构成素谱的子集．

\begin{enumerate}
    \item \textbf{空集与全空间}：\(V(1) = \varnothing\)，因为不存在包含单位理想的素理想；\(V(0) = \operatorname{Spec} R\)，因为零理想包含于任何素理想．
    \item \textbf{有限并封闭}：对任意理想 \(I, J\)，有 \(V(I) \cup V(J) = V(I \cap J)\)．事实上，\(V(I \cap J) \supseteq V(I) \cup V(J)\) 显然；反之，若 \(\mathfrak{p} \supseteq I \cap J\) 且 \(\mathfrak{p}\) 不含 \(I\) 也不含 \(J\)，则存在 \(x \in I \setminus \mathfrak{p}, y \in J \setminus \mathfrak{p}\)，则 \(xy \in I \cap J \subseteq \mathfrak{p}\)，由素理想性得 \(x \in \mathfrak{p}\) 或 \(y \in \mathfrak{p}\)，矛盾．故 \(V(I \cap J) \subseteq V(I) \cup V(J)\)．更常用的是 \(V(I) \cup V(J) = V(IJ)\)，因为 \(IJ \subseteq I \cap J\)，且 \(V(IJ) = V(I \cap J)\)．
    \item \textbf{任意交封闭}：对任意一族理想 \(\{I_\alpha\}\)，有 \(\bigcap_\alpha V(I_\alpha) = V(\sum_\alpha I_\alpha)\)，因为 \(\mathfrak{p}\) 包含所有 \(I_\alpha\) 当且仅当它包含它们的和．
\end{enumerate}

以退化集为闭集，
使得素谱$\text{Spec} A$成为拓扑空间，
称为\textbf{Zariski 拓扑}．
对于多项式环，
它对应着零点闭集的Zariski拓扑．

对于 $f \in A$，
不包含它的素理想构成
$V((f)) = \{ \mathfrak{p} \mid f \notin \mathfrak{p} \}$．

定义 $D_f = \text{Spec} A \setminus V((f))$．
这些 \(D(f)\) 构成 Zariski 拓扑的一族开基，称为\textbf{基本开集}．它们在代数几何中对应非零函数的非零点集．


前述的扩环产生了自然的
$\mathbb Z \to \mathbb Z[\sqrt d]$的嵌入，
使$\mathbb Z$成为子环．
扩环的过程改变了可约性，
因此改变了素理想和素谱．
Fermat在Diophantus著作上批注了命题：
若有素数$p \equiv 1 \pmod{4}$ iff.
$p$可以表示为整数平方和$p = x^2 + y^2$．
用现代的观点看，
在$\mathbb Z$中，
模$4$求余得到$\mathbb Z/(4) \simeq {0,1,2,3}$，
余数为偶数的均不是素数，
只剩下余数为$1,3$的情况．
当素数$p$的余数为$3$时，
素数$p$在$\mathbb Z[\text i]$中仍然为素数，
即不可分解的Gauss素数．
最后一种情况即
$p \equiv 1 \pmod{4}$ iff. $(p - 1) \in (4)$，
在$\mathbb Z[\text i]$中，
有不可约的分解：
$p - 1 = x^2 + y^2 - 1 = (x + y\text i)(x - y\text i) - 1 \in (4)$，
例如
$5 - 1 = (2 + \text i)(2 - \text i) - 1$，
$13 - 1 = (3 + 2\text i)(3 - 2\text i) - 1$，
$17 - 1 = (4 + \text i)(4 - \text i) - 1$

整数环 \(\mathbb{Z}\) 的素谱为
$\text{Spec} \mathbb Z = \{ (0) \} \cup \{ (p) \mid p \text{ 是素数} \}$，
其中 $(0)$ 是泛点（generic point），
每个素数 $p$ 对应一个闭点 $(p)$．
Zariski 拓扑下，
闭集由有限多个素理想组成（以及整个空间）．
一个理想 $I = (n)$ 的退化集 $V(I)$ 包含所有整除 $n$ 的素数生成的理想．

\section{剩余域}
  
\begin{framed}
\noindent
交换环 $A$ 对其极大理想 $\mathfrak m$ 构造的商环 $A/\mathfrak m$ 是域，
称为\textbf{剩余域}（Residue Field）．
\end{framed}

若整环$A$对其素理想$\mathfrak p\in\text{Spec} \ A$局部化得到 $A_\mathfrak p$ 和唯一的极大理想 $\mathfrak p A_\mathfrak p$，
则可以对这个极大理想做剩余域
$k(\mathfrak p) = A_\mathfrak p/\mathfrak p A_\mathfrak p$．
此外商环$A/\mathfrak p$是整环，
有分式域$\text{Frac} \ A/\mathfrak p = (A/\mathfrak p\setminus \{0\})^{-1}(A/\mathfrak p) $．
有典范同构
$k(\mathfrak p) \simeq \text{Frac} \ A/\mathfrak p$．

\begin{equation}
\mathfrak p 
\mapsto (S = A \setminus \mathfrak p)
\mapsto (A_\mathfrak p = S^{-1} A)
\mapsto \mathfrak pA_\mathfrak p
\mapsto \big( k(\mathfrak p) = A_\mathfrak p/\mathfrak p A_\mathfrak p \simeq \text{Frac} \ A/\mathfrak p \big)
\end{equation}
\label{eq:prime-ideal-local-ring}

对于极大理想$\mathfrak m$，
$A/\mathfrak m$已经是剩余域，
$\text{Frac} \ A/\mathfrak m = A/\mathfrak m$．

\begin{equation}
\mathfrak m
\mapsto \mathfrak m A_\mathfrak m
\mapsto \big( k(\mathfrak m) = A_\mathfrak m/\mathfrak m A_\mathfrak m \simeq \text{Frac} \ A/\mathfrak m = A/\mathfrak m \big)
\end{equation}
\label{eq:max-ideal-local-ring}

$3\mathbb{Z}_{(3)} = \{ 3a/s  \mid a \in \mathbb Z,\ 3 \nmid s \}$ 作为不可逆元构成的集合，
成为了$\mathbb{Z}_{(3)}$的唯一极大理想，
剩余域为
$k((3)) = \mathbb{Z}_{(3)}/(3)\mathbb{Z}_{(3)}$，
它同构于$\text{Frac} \mathbb Z/(3)$，
即Galois域$\mathbb F_3$．


\begin{itemize}
    \item 对局部环 $\mathbb{Z}_{(p)}$ 模掉它的极大理想 $p\mathbb{Z}_{(p)}$，得到的剩余域同构于 Galois 域 $\mathbb{F}_p$（即 $\mathbb{Z}/p\mathbb{Z}$）．
    \item 在实系数多项式环 $\mathbb{R}[x]$ 中，对极大理想 $(x^2+1)$ 进行商运算，其剩余域即为复数域 $\mathbb{C}$．
\end{itemize}
在微分几何与代数几何的统一视角中，剩余域不仅能被用来提取“点”的坐标值，还能通过极大理想的平方 $\mathfrak{m}^2$ 构造切空间．商空间 $\mathfrak{m}/\mathfrak{m}^2$ 模除了所有的二阶高阶小量，天然同构于该点处的余切空间．

$(x - 1) \in \text{Spec} \ \mathbb R[x]$．
$(x^2 + 1) \notin \text{Spec} \ \mathbb R[x]$，

\begin{equation}
\begin{aligned}
0 \to (x - 1) \to \mathbb R[x] \to \mathbb R \to 0 \\
0 \to (x^2 + 1) \to \mathbb R[x] \to \mathbb C \to 0 \\
\end{aligned}
\end{equation}
\label{eq:short-exact-seq-num-field}

即$\mathbb R[x]/(x - 1) \simeq \mathbb R$以及$\mathbb R[x]/(x^2 + 1) \simeq \mathbb R$．
更多阐述见(\ref{eq:short-exact-seq-fun-real})
\chapter{层}

拓扑空间$X \in\textbf{Top}$上构造交换环层$\mathcal O_X$，
点$x \in X$上有局部环$\mathcal O_x$，
$\exists! \mathfrak m_x$，
构成局部赋环空间$(X,\mathcal O_X)$

实拓扑流形$X$上构造实连续函数层$\mathcal O_X$，
剩余域为实域．
实光滑流形$X$上构造实光滑函数层$\mathcal O_X$，
剩余域为实域．
复流形$X$上构造全纯函数层$\mathcal O_X$，
剩余域为复域．
概形$X$上构造正则函数层$\mathcal O_X$．

层范畴$\textbf{Sh}(X)$是Grothendieck topos，
$\mathcal O_X$视为其中的环对象．


\section{完备格}

\begin{framed}

$(L,\leq) \in\textbf{Poset}$是偏序集．
若对于$\forall a,b \in L$，
二元上确界(并join)$a\vee b$
和下确界(交meet)$a\wedge b$都存在，
即$\vee,\wedge$运算封闭，
称$L$为格．

若存在$0,1 \in L$ s.t.
$\forall x \in L$有$0 \leq x \leq 1$，
即$x \vee 0 = x,\ x \wedge 1 = x$，
或者等价的$x \vee 1 = 1,\ x \wedge 0 = 0$，
称$L$为有界格，
$0$为下确界，
$1$为上确界．

若任意子集$S \subseteq L$都有上确界
$\text{sup} \ S = \bigvee S$
和下确界
$\text{inf} \ S = \bigwedge S$，
称$L$为完备格．

\end{framed}

\begin{example}{$PX$}

幂集$PX$按照包含关系构成有界格和完备格，
格的交和并运算对应集合的交和并运算，
格有下确界$\varnothing$，
上确界$X$．

\end{example}

\begin{example}{$[0,1]$}

闭区间$[0,1]$按照自然序构成有界格和完备格，
格的交和并运算对应数的下确界和上确界运算．
格有下确界$0$，
上确界$1$．

\end{example}

\begin{example}{$\mathbb Z$}

$\mathbb Z$按照整除性$a \mid b$构成格，
$a \vee b = \text{lcm}(a,b)$，
$a \wedge b = \text{gcd}(a,b)$．
整数的无限子集，
例如素数的全体，
没有最小公倍数，
故$\mathbb Z$不是有界格，
更不是完备格．

若在理想上而非集合上考虑，
令$\mathcal I_\mathbb Z$为$\mathbb Z$的理想的全体．
有$a \mid b$ iff. $(b) \subseteq (a)$，
即便全体素数构成的无限子集$\{p_i\}$没有最小公倍数，
其理想$(p_i)$的交为$(0) \in \mathcal I_\mathbb Z$，
故下确界始终存在．
故$\mathcal I_\mathbb Z$是完备格．

\end{example}

交换环$A\in\textbf{CRing}$的理想全体
$\mathcal I_A$是有界格也是完备格．
以上例子告诉我们，
有时候在理想构成的完备格上，
相比环子集构成的格具有更好的性质．

\begin{framed}
格$(L,\vee,\wedge)$若满足以下相互等价的等式：
左分配律$a\wedge(b\vee c) = (a\wedge b) \vee (a\wedge c)$，
右分配律$a\vee(b\wedge c) = (a\vee b) \wedge (a\vee c)$，
则具有分配律．
\end{framed}

最常见的例子即是子集运算的分配律．

\section{Heyting代数}

\begin{framed}
有界分配格$(H,\vee,\wedge,0,1)$，
若有称之为蕴涵的二元封闭运算$\Rightarrow$，
$\forall a,b \in H$，
$\exists (a \Rightarrow b) \in H$
s.t. 
$\forall c \in H$，

$$
c \leq (a \Rightarrow b)
\Longleftrightarrow
(a \wedge c) \leq b
$$

称$(H,\vee,\wedge,\Rightarrow,0,1)$为Heyting代数．

定义补运算为：

$$
\neg a = (a \Rightarrow 0)
$$

\end{framed}

以上定义中的条件写为偏序范畴的态射集
$H(c,a \Rightarrow b) \simeq H(a \wedge c,b)$，
以便发现这里实际上是一对伴随函子$(a \wedge -) \dashv (a \Rightarrow -)$：

$$
% https://tikzcd.yichuanshen.de/#N4Igdg9gJgpgziAXAbVABwnAlgFyxMJZABgBoAmAXVJADcBDAGwFcYkQAKACVIB1fGMAI4BKEAF9S6TLnyEUZYtTpNW7bnwHCxk6djwEiAFgrKGLNohAAjCVJAZ9c46SU1zaq-QAE-AO4wUADmMN4AxnZ6soYo5K5mqpacPPyCopEOMgbyyHFU7onqKVrpuplOMcgAbPEFFuwRZY7ROTX5KvVevrwASlhBABY49ABOIxB+3rbiyoEhCCigAGbjALZIAIw0OBBIZCBwA1hLOHt1niCpwhkrEOuIAMzbu4hxHRdXQiA0jPTWMIwAApZZxWEb9IY3NZIEwgHZIACs5ySnyhdxhz02yPYPn8c1CAFpviBfv8gSCYiBwYNTmVbvd9vDEEj3klcb0IcMxhNvESfn8AcCKvIqZy0fcanCXk9WTjugFgoTiaTBRSRdTIXToYhJUyAOzYqyorXo16YxAG2Vdfh9GmjcaTPkkgXk4XsDW0yjiIA
\begin{tikzcd}
{(H,\leq)} \arrow[dd, "\leq"]             &  & {(H,\leq)} \arrow[dd, "\leq"] \arrow[ll, "a \wedge -"'] &  & a \wedge c \arrow[dd, "\leq"']   &  & c \arrow[ll, "a \wedge -"'] \arrow[dd, "\leq"] \\
                                          &  &                                                         &  &                                  &  &                                                \\
{(H,\leq)} \arrow[rr, "a \Rightarrow -"'] &  & {(H,\leq)}                                              &  & b \arrow[rr, "a \Rightarrow -"'] &  & a \Rightarrow b                               
\end{tikzcd}
$$

以上的蕴涵运算$a \Rightarrow b$常令人疑惑．
在经典逻辑中，
蕴涵写为$A \to B$，
直觉上理解为：
“若A真，则B真”．
这个陈述为假则意味着：
“A真而B假”即$A \wedge \neg B$，
除此之外都应该视为真：

$$
\neg (A \wedge \neg B) = \neg A \vee B
$$

这里的$\rightarrow$和$\neg$都符合Heyting代数的定义，
为防止二者的循环引用，
可以直接用真值表来定义蕴涵．

\begin{example}{$PX$}

幂集$PX$上，
有补集运算，
以及参照逻辑的方式定义蕴涵
$(A \Rightarrow -) = A^c \cup -$，
有
$C \subseteq A^C \cup B 
\Longleftrightarrow 
A \cap C \subseteq B$．

\end{example}

\begin{example}{$\text{Open} \ X$}

拓扑空间$X$上的开集范畴$\text{Open} \ X$
是$X$的子集族．
和$PX$对补的定义不同，
开集的补集不一定是开集，
为保持补运算的封闭性，
定义开集$A$的补$A^c$为$X \setminus A$中最大的开集，
再按照$PX$相似的方式定义蕴涵
$(A \Rightarrow -) = A^c \cup -$，

$C \subseteq A^C \cup B 
\Longleftrightarrow 
A \cap C \subseteq B$．

\end{example}

Heyting代数上自带的伴随函子，
使得它构成CCC．
蕴涵$\Rightarrow$可以构造指数对象．


\section{生成理想}

理想作为子集可以覆盖给定的子集$S$，
而理想又可以被素理想所覆盖，
这样就引入了子集$S$所生成的理想，
以及相应的素理想问题．

\begin{example}{$\mathbb Z$}

子集$S=\{ -60, 120, 720 \}$，
覆盖它的理想集合为$\{(1),(2),(3),(5),(6),(15),(60)\} \in I_\mathbb Z$，
最大公约数$\text{gcd} \ S = 60$，
理想$(\text{gcd}) = (60) = (1)\cap(2)\cap(3)\cap(5)\cap(6)\cap(15)\cap(60)$．

$$
\begin{tikzcd}
     &  &                  &  & (2) \arrow[lld] \arrow[lllldd, bend right] \\
     &  & (6) \arrow[lld]  &  &                                            \\
(60) &  &                  &  & (3) \arrow[llu] \arrow[lld] \arrow[llll]   \\
     &  & (15) \arrow[llu] &  &                                            \\
     &  &                  &  & (5) \arrow[llu] \arrow[lllluu, bend left] 
\end{tikzcd}
$$

按照理想升链，
跳过中间过程的理想$(6),(15)$，
$S$最终被极大素理想$(2),(3),(5)$所覆盖，
记

$$
V(S) = \{(2),(3),(5)\}
$$

取整系数$\lambda_i\in \mathbb Z$做有限形式和
$\lambda_1 \cdot (-60) + \lambda_2 \cdot 120 + \lambda_3 \cdot 720$，
其全体构成理想$(S) = \{ \lambda_1 \cdot (-60) + \lambda_2 \cdot 120 + \lambda_3 \cdot 720 \}$．

\end{example}

\begin{example}{$\mathbb R[X]$}

$X = \mathbb R$，
多项式环$\mathbb R[X]$的子集
$S=\{ f_1(x) = x - a, f_2(x) = x - b, f_3(x) = x-3 \}$，
有零点$Z(S) = \{ a, b, c\}$．
$S$的任何子集$S_i$的零点集
$Z(S_i)$都覆盖了$S$，
这些零点集的交为：
$\bigcap\limits_i Z(S_i)$

$$
% https://tikzcd.yichuanshen.de/#N4Igdg9gJgpgziAXAbVABwnAlgFyxMJZAFgBoAGAXVJADcBDAGwFcYkQAtACnoEoQAvqXSZc+QigBMpAIzU6TVu271SAI35CR2PASLlSk+QxZtEnHutIBjTcJAYd4omSM0TS89w2D7jsXpSpADMxopmFmo2dtoBEiSkxGGmyly2gvIwUADm8ESgAGYAThAAtkgGIDgQSDLu4UhgzIyMNIz0ajCMAAqiuhIgjDAFOL6FJeWIldVI0gopiE0tbR1dvU6Bg8OjNJ1gUEgAtMHkWiDFZbU0M4hzHmZLrYOrPX3O5kMjY+cTSMHXNUQdXmnkeK06rw2A0+ozOF0m-yqgLIIIezSe7Qh6zi7Bh33hfwBs3qCzBzyxb02eLhv0QKJudwai3R4LWlOh23xtIArES6STQSzyWyobjOTTLoheUjiajGkLMSKcR9ObsYPs-qdKAIgA
\begin{tikzcd}
           &  &                      &  & Z(a) \arrow[lld] \arrow[lllldd, bend right] \\
           &  & {Z(a,b)} \arrow[lld] &  &                                             \\
{Z(a,b,c)} &  &                      &  & Z(b) \arrow[llu] \arrow[lld] \arrow[llll]   \\
           &  & {Z(b,c)} \arrow[llu] &  &                                             \\
           &  &                      &  & Z(c) \arrow[llu] \arrow[lllluu, bend left] 
\end{tikzcd}
$$

取环系数$\forall \lambda_i \in A$
做有限和$\sum_i \lambda_i f_i \in A$，
其全体构成理想
$(S) = \{ \sum_i \lambda_i f_i \in A \}$．

\end{example}

子集$S$生成的理想$S$有如下两种定义：

\begin{framed}
\noindent
$(S) = \bigcap\limits_{I \supset S,\ I \in I_A} I$
\end{framed}

\begin{framed}
\noindent
$(S) = \left\{ \sum_{i=1}^n a_is_i \mid n \in \mathbb N,\ (a_i,s_i) \in A \times S \right\}$
\end{framed}

这样可以引入主理想和PID的概念：

\begin{framed}
\noindent
由单个元素 $a \in A \in \textbf{CRing}$ 生成的理想 $(a) = \{ra \mid r \in A\}$ 称为\textbf{主理想}（Principal Ideal）．
若整环的所有理想都是主理想，则称其为\textbf{主理想整环}（PID）．
\end{framed}

$\mathbb{Z}$ 是 PID 的典型例子；而多项式环 $\mathbb{Z}[x]$ 并非 PID（例如理想 $(2, x)$ 不能由单元素生成）．

\begin{framed}
\noindent
由单个元素 $a \in A \in \textbf{CRing}$ 生成的理想 $(a) = \{ra \mid r \in A\}$ 称为\textbf{主理想}（Principal Ideal）．
若整环的所有理想都是主理想，则称其为\textbf{主理想整环}（PID）．
\end{framed}

\section{反单调Galois连接}

$A$的幂集有自然的偏序结构$(\textbf 2^A,\leq)$，
把子集$S \in \textbf 2^A$置于这个偏序结构考虑．
对函数环$A = C(X)$的子集$S$求零点集，
视为幂集上的函数：
$Z: \textbf 2^A \to \textbf 2^X:: S \mapsto Z(S)$．
反过来，
子集$E \in X$确定了理想
$I(E) = \{f \in A \mid \forall x \in E:\ f(x) = 0\}$，
视为幂集上的函数：
$I: \textbf 2^X \to \textbf 2^A :: E \mapsto I(E)$．

前面讨论过，
$\textbf 2^X,\textbf 2^A$是偏序集和完备格，
但我们需要讨论的是$A$中理想之间的结构，
而非$A$的子集之间的结构，
需要转换到$I_A$这个完备格上讨论．
将子集$S\in\textbf 2^A$用$(S)\in \text{Ideal} \ A$来替换，
得到
$Z: I_A \to \textbf 2^X:: (S) \mapsto Z((S)) = Z(S)$，
以及
$I: \textbf 2^X \to I_A :: E \mapsto I(E)$．

\begin{framed}
\noindent

反单调Galois连接是连接两个偏序范畴$(P,\leq)$和$(Q,\leq)$的反变的伴随函子，
$q \leq Lp$ iff. $p \leq Rp$对应了
$(Q,\leq)(q,Lp) \simeq (P,\leq)(p,Rq)$：

$$
% https://tikzcd.yichuanshen.de/#N4Igdg9gJgpgziAXAbVABwnAlgFyxMJZABgBoAmAXVJADcBDAGwFcYkQAKARVIB1fGMAI4BKAHr8cMAB45gENAF8Qi0uky58hFGWLU6TVu258Bw8ZJlyFy1eux4CRACwV9DFm0QghKtSAwHLRdSPRoPI28AGTQ-e00nFHJQ90MvTgAFU0FROICNR21kZKpwtOMs-hyRPMCEooA2FLLPdli7fKDE5CbSg1bvACVfRX0YKABzeCJQADMAJwgAWyQyEBwIJABGGjgACyxZnFWWyJAq4TyF5aRk9c3EAGZT9IvfGkZ6ACMYRgyC4LeQRHK6LFaIVz3JAAVhe7DeoJuEJoG22cOiIA+31+-y62hA8ywEz2xw613Ba1RiFh-TOg0xIE+Pz+AMSBKJJMR4KaUKe6JAUQZTJxrPxhOJpP85KQPKpAHZ+QiyWDbiiHgrael6Vjmbj6uxxZzRoogA
\begin{tikzcd}
{(Q,\leq)^\text{op}}                                     &  & {(P,\leq)} \arrow[dd, "\leq"] \arrow[ll, "L"'] &  & Lp                                    &  & p \arrow[ll, "L"'] \arrow[dd, "\leq"] \\
                                                         &  &                                                &  &                                       &  &                                       \\
{(Q,\leq)^\text{op}} \arrow[uu, "\leq"] \arrow[rr, "R"'] &  & {(P,\leq)}                                     &  & q \arrow[uu, "\leq"] \arrow[rr, "R"'] &  & Rq                                   
\end{tikzcd}
$$

\end{framed}

这样$(Z,I)$构成反单调Galois连接：

$$
% https://tikzcd.yichuanshen.de/#N4Igdg9gJgpgziAXAbVABwnAlgFyxMJZABgBoAmAXVJADcBDAGwFcYkQAKAHS5xgA8cAAnIA9ABqkejGAEcAlKJ59BwCGgC+IDaXSZc+QijLFqdJq3bdeA4WMnS5i5bbWbtukBmx4CRACwUZgwsbIggAKIeej6GAaSmNCGW4QBaHBwAyvLy0V76vkbI5AnBFmGcAJIA+gCCUlwyCnneBn4oJVRJ5VY19Y7NOjFtRQBspd2hVtktBXEo413mU+GVHBG5GmYwUADm8ESgAGYAThAAtkhkIDgQSACMNHAAFlhHOFeTKSADeacXSBKNzuiAAzF8Kr8aIx6AAjGCMAAKc3aIBk7z+Z0uiECwKQAFYIexfkMQP9sbjbg8iWkQNC4QjkbFUScsLtnh9SeTPnjEITlt9KnS0QykSijCBWezOZ5uYhxrzwQKKqlhTD4WLmRKpRzMQD5TQqYgAOw0n6NOR67FAo2m5XsIX0jVMkbsHWcygaIA
\begin{tikzcd}
{(\text 2^X,\leq)^\text{op}}                                     &  & {(I_A,\leq)} \arrow[dd, "\leq"] \arrow[ll, "Z"'] &  & Z((S))                                &  & (S) \arrow[ll, "Z"'] \arrow[dd, "\leq"] \\
                                                                 &  &                                                  &  &                                       &  &                                         \\
{(\text 2^X,\leq)^\text{op}} \arrow[uu, "\leq"] \arrow[rr, "I"'] &  & {(I_A,\leq)}                                     &  & E \arrow[uu, "\leq"] \arrow[rr, "I"'] &  & I(E)                                   
\end{tikzcd}
$$

\section{函数环}

实微分几何的基本研究对象是实光滑流形$X$上的实光滑函数环$C(X)$．
在$X$上选定子集$S$得到零点集 
$Z: \textbf 2^X \to \textbf 2^{C(X)} :: S \mapsto Z(S) = \{ f \in C(X) \mid \forall x \in S,\ f(x) = 0 \}$．
这构造了幂集之间的反单调偏序Galois连接．
除非$X$是单点，否则$C(X)$不是整环．
    
每个点$x \in X$都可以构造退化函数集合作为极大理想
$X \to \text{Max} \ C(X) : x \mapsto \mathfrak m_x = Z(\{ x \}) = \{ f \in C(X) \mid f(x) = 0 \}$．
构造$x$点上的赋值同态：

$$
\begin{aligned}
\varphi: C(X)/\mathfrak m_x &\to \mathbb R \\
f + \mathfrak m_x &\mapsto \varphi(f + \mathfrak m_x) = f(x) \\
\end{aligned}
$$

常值函数可以取值任意实数，
故$\varphi$是满的．
同态的核$\text{Ker} \ \varphi = \{ 0 \}$，
故$\varphi$是单的，
故$\varphi$是同构即

\begin{equation}
C(X)/\mathfrak m_x \simeq \mathbb R
\end{equation}
\label{eq:fun-ring-residue-field-real}

进而构造乘法系统$C(X)\setminus \mathfrak m_x$，
局部化：
$\mathcal O_x = C(X)_{\mathfrak{m}_x}$，
其元素为分式 $f/g$，
其中 $(f,g \in C(X),\ g(x) \neq 0$．
$f/g = f'/g'$ iff. $\exists h \neq 0$ s.t. $h(f g' - f' g) = 0$．
即在点 $x$ 的某个邻域上 $f/g = f'/g'$．
因此 $\mathcal O_x$ 恰好同构于 $x$ 点的光滑函数芽环，
即所有在 $x$ 的某个邻域上有定义的光滑函数在“局部相等”关系下的等价类构成的环．
根据(\ref{eq:fun-ring-residue-field-real})，
类似(\ref{eq:short-exact-seq-num-field})：

\begin{equation}
0 \to \mathfrak m \to \big( \mathcal O_x = C(X)_{\mathfrak{m}_x} \big) \to \mathbb R \to 0
\end{equation}
\label{eq:short-exact-seq-fun-real}


\section{函数层}

拓扑空间$X$的开集族$\text{Open} \ X$在嵌入包含下构成偏序范畴，
函数层的构造是反变的偏序函子$\mathcal O$：

$$
% https://tikzcd.yichuanshen.de/#N4Igdg9gJgpgziAXAbVABwnAlgFyxMJZABgBpiBdUkANwEMAbAVxiRAAoAdTnGADxzAA8mhhgAvgAJukgBqlpnOEwBGcGLwCOASgB63XgOAQ04kONLpMufIRQAmclVqMWbLj344VAM2ABhACUsMABzC0VDQQAneHFtc0sQDGw8AiIyAEZnemZWRBAAVUSrVNsiR2zqXLcC7gBbOhwACwBjRkkhAH1CyQBeSX9ii1KbdJQyAGYc13yQADUS5Os0u2RHaerZtgamto7u+f7BxfFnGChQ+CJQH2iIeqQyEBwIJEytvJ3ORpb2hk6XVkSzuDyQjheb0QkxGIFBj0QABZqK8kABWWHw8EoqHIlxfOpKVTqLQgagMOgqGAMAAKK3KBWiWFCzRwIPuCLROKQk0+tRABi8wFicDM5Mp1LpZXGICZLLZZ3EQA
\begin{tikzcd}
{(\text{Open} \ X, \subseteq)^\text{op}} \arrow[rr, "\mathcal O_X"] &  & {(\textbf{CRing}, \text{res})}              \\
U \arrow[rr] \arrow[dd, "\subseteq"']                               &  & \mathcal O_U = CU                           \\
                                                                    &  &                                             \\
V \arrow[rr]                                                        &  & \mathcal O_V = CV \arrow[uu, "\text{res}"']
\end{tikzcd}
$$

限制映射$\text{res}: \mathcal O_V \to \mathcal O_U :: f \mapsto f\mid_U$．
函子$\mathcal O_X$是$X$上的 \textbf{结构层} (structure sheaf)．
这样的函子的全体构成函数层范畴$\textbf{Sh}(\textbf{CRing})$．

为了考虑$X$上的局部性质，
需要让局部以外的性质被忽略等价商掉．
比较函数$f,g\in\mathcal O_X = C(X)$，
其差异$f-g$的全局性质，
随着限制的极限过程而消去：
$(f-g)\mid_{U_1} \stackrel{\text{res}}{\longrightarrow}(f-g)\mid_{U_2} \stackrel{\text{res}}{\longrightarrow} \cdots $，
在这一极限过程中，
若存在点$x$的开邻域$U$ s.t. $(f-g)\mid_U = 0$，
这意味着$f,g$在$x$点的局部是相等的，
构成等价关系$f \sim_x g$，
等价类称为函数芽$[f]_x = [g]_x$，
芽构成的环称为\textbf{茎}(stalk)：
$\mathcal O_x = \varinjlim_{U \ni p} \mathcal O_U$．

在$x$处退化的函数芽构成极大理想
$\mathfrak m_x = \{ [f] \in \mathcal{O}_x \mid f(x) = 0 \}$．
由于茎的构造基于拓扑空间$X\in\textbf{Top}$，
函数的连续性是基本要求，
若不加说明函数环都是连续函数环．
在极大理想之外的芽$f \notin \mathfrak m_x$是非退化的，
由于函数的连续性保证了$f$的可逆性．

这是唯一的极大理想，
使得茎$\mathcal O_x$是局部环，
以及 $(X, \mathcal O_X)$ 构成一个\textbf{局部赋环空间}（locally ringed space）．



\chapter{同调}

\section{自由Abel群}

Abel群$A$中的元素$x$有逆元$-x$，
整数$n \in \mathbb Z$给出了

$$
nx = \begin{cases}
  \underbrace{x + \cdots + x}_n & (n > 0) \\
  0 & (n = 0) \\
  \underbrace{(-x) + \cdots + (-x)}_{|n|} & (n > 0) \\
\end{cases}
$$

进而取子集$\{x_i\} \subseteq A$构造和：

$$
\sum_i n_i x_i
$$

这样形式的全体构成了$\{x_i\}$生成的子群。
若生成集合有限则称有限生成群，
若$\sum_i n_i x_i = 0$的条件是所有$n_i = 0$退化，
称生成集$\{x_i\}$线性无关。
由$r$个线性无关元有限生成的群称为秩$r$自由Abel群。
$\mathbb Z$是$1$或$-1$生成的自由Abel群。


\section{正合}

环$\Lambda \in \textbf{Ring}$上的左模范畴记$\Lambda-\textbf{Mod}$，
考虑模同态$f \in \Lambda-\textbf{Mod}(X,Y)$，
产生了同态的像与核：

$$
\text{Im} \ f = \{ y \in Y \mid \exists x \in X,\ f(x) = y\},\ 
\quad
\text{Ker} \ f = \{x\in X \mid f(x) < 0\}
$$

它们都是子模，
从而可以用商模定义：

$$
\text{Coim} \ f = X/\text{Ker} \ f,\ 
\quad
\text{Coker} \ f = Y/\text{Im} \ f
$$

由$f$有同态：

$$
\begin{aligned}
  f^\prime : \text{Coim} \ f = X/\text{Ker} \ f &\to Y \\
  x + \text{Ker} \ f &\mapsto f^\prime(x + \text{Ker} \ f) = f(x)
\end{aligned}
$$

显然$\text{Im} \ f^\prime = \text{Im} \ f$，
$f^\prime$是满同态。
$f^\prime(x + \text{Ker} \ f) = f(x) = 0$时，
$x \in \text{Ker} \ f$，
即$x + \text{Ker} \ f = 0$，
有$\text{Ker} \ f^\prime = 0$，
$f^\prime$是单同态。
$\text{Coim} \ f$描述的就是同构$ = X/\text{Ker} \ f \simeq \text{Im} \ f$。
当$\text{Ker} \ f = 0$时称为单同态，
$\text{Im} \ f = Y$时称为满同态，
单满同态为同构，
此时

$$
\text{Coim} \ f = X/\text{Ker} \ f \simeq X \simeq \text{Im} \ f = Y
$$

像和核都是子模，
若将二者放在一个模中比较，
即两个首尾相接的同态中，
$\text{Im} \ f$和$\text{Ker} \ g$都是$Y$的子模：

$$
X \stackrel{f}{\longrightarrow} Y \stackrel{g}{\longrightarrow} Z
$$

若复合映射$X \stackrel{g \circ f}{\longrightarrow} Z$
退化为$g \circ f = 0$，
此时$\text{Im} \ f \subseteq \text{Ker} \ g$。
当$\text{Im} \ f = \text{Ker} \ g$时称为正合。
正合可以用于描述单/满同态：

$$
0 \stackrel{i}{\longrightarrow} X \stackrel{f}{\longrightarrow} Y,\
\quad
Y \stackrel{g}{\longrightarrow} Z \stackrel{p}{\longrightarrow} 0
$$

以上两个正合意味着$\text{Im} \ i = \text{Ker} \ f = 0$，
它说明$f$是单同态；
以及$\text{Im} \ g = \text{Ker} \ p = Z$，
它说明$g$是满同态。
将二者合并得到短正合序列，
描述了单同态$f$和满同态$g$在$Y$处正合：

$$
0 \longrightarrow X \stackrel{f}{\longrightarrow} Y \stackrel{g}{\longrightarrow} Z \longrightarrow 0
$$

单同态$f$把$X$以同构的意义嵌入$Y$作为子模，
在自然投射$\pi: Y \to Y/X$中，
子模$X \simeq \text{Ker} \ \pi$，
于是：

$$
0 \longrightarrow X \stackrel{f}{\longrightarrow} Y \stackrel{\pi}{\longrightarrow} Y/X \longrightarrow 0
$$


\section{复形}


\section{同调群}




\chapter{流形}

\section{函数芽}

局部化环$C(X)_{\mathfrak m_x} = \{ f/g \mid f,g \in C(X),\ g(x) \neq 0 \}$，
茎环$\mathcal O_x = \varinjlim_{U \ni p \ni x} \mathcal O_U$．
微分几何研究的光滑流形$X$是仿紧的Hausdorff空间，
其拓扑性质保证了局部化环的同构：

$$
\mathcal O_x \simeq C(X)_{\mathfrak m_x}
$$

局部化环$A_\mathfrak p$的唯一极大理想是$\mathfrak p A_\mathfrak p$，
这里同构的局部化环的唯一极大理想就是
$\mathfrak m_x \mathcal O_x$．
根据(\ref{eq:max-ideal-local-ring})和(\ref{eq:fun-ring-residue-field-real})有
$\mathcal O_x/\mathfrak m_x \mathcal O_x \simeq C(X)/\mathfrak m_x \simeq \mathbb R$．
这反映了在点 $x$ 附近函数的值由其芽唯一确定．


实光滑流形的光滑函数环$C(X)$在点$x$上有极大理想
$\mathfrak m_x = \{ f \in C(X) \mid f(x) = 0 \}$
和局部环$\mathcal O_x \simeq C(X)_{\mathfrak m_x}$．
用$\mathfrak m_x$的元素有限生成理想
$\mathfrak m_x^2 = \{ \sum_i f_i g_i \mid f_i,g_i \in \mathfrak m_x \}$．
其意义在于，
$\mathfrak m_x$中的函数在$x$处退化（一阶零点），
$\mathfrak m_x^2 \subset \mathfrak m_x$中的函数在$x$处退化，
且其微分也退化（二阶零点）．
余切空间在同构意义下等于
$ T^*_x X \simeq \mathfrak m_x/\mathfrak m_x^2 $，
其元素记为：

$$
\text d f = f + \mathfrak m_x^2 \in T^*_x X
$$

例如，
$X = \mathbb R$，
$x = 0$，
而$f(x) = 3x \in C(X)$在$x$处只具有一阶零点，
故有非零的等价类：
$\text d f = f + \mathfrak m_x^2$．
$f(x) = x^2 \in C(X)$在$x$处具有二阶零点，
即$f \in \mathfrak m_x^2$，
故等价类退化：
$\text d f = 0$．

二维情形$X = \mathbb R^2$，
$(x,y) = (0,0)$，
而$f(x) = 3x + 2y \in C(X)$在$(x,y)$处只具有一阶零点，
故有非零的等价类：
$\text d f = f + \mathfrak m_{(x,y)}^2$．
$f(x) = 3x^2 + 2y^2 \in C(X)$在$(x,y)$处具有二阶零点，
即$f \in \mathfrak m_{(x,y)}^2$，
故等价类退化：
$\text d f = 0$．


\section{对偶空间}

            目前教育界在线性代数的教学上缺少一些重要的概念——对偶空间．对偶是现代数学最核心的数学思想之一．最简单的对偶是体现在线性代数中线性向量和线性函数之间的相互对偶关系之中的．\par
            本文是在一般线性代数的基础上，介绍高级一些的线性代数知识，以便将来在学习微分几何和泛函分析时，对对偶、逆变、协变这些概念有基本的理解．内容概要：\par
            介绍线性空间的线性映射、对偶空间，行向量和列向量的对偶表示，双线性映射，Einstein求和约定．\par
            接着介绍基的变换，以及在这种变换下，对偶空间上的向量和对偶向量的变换规则，即逆变和协变规则．\par
            继续以平面旋转群SO(2)为例子，研究在基进行平面旋转的情况下，向量的变换规律．\par
            介绍对偶基的概念．一对对偶的向量是通过内积的结构成为双线性映射的．略微介绍Riesz引理，了解到对偶的向量之间如何一一对应，从而让内积可以直接在逆变向量之间实现．\par
            前面几讲中用了许多 $\mathbb{R}^2$ 中的技术来讨论复平面的问题，体现出线性代数基础的重要性．这一讲要谈线性代数最核心的对偶性问题，在对偶性上展开讨论一些重要的基础概念．我们在线性代数上有了这些基础之后，反过来又可以了解复结构是如何应用于对偶性空间的．\par
            线性代数学习中，我们的讨论多专注在 $\mathbb{R}^n$ 上．若能理解概念，在一般的线性空间（有限/无限维）上扩展概念并非难事，参考一般的线性代数/泛函分析教材即可．\par
            在另一讲中用到一个非常好的例子阐述对偶的基本原理，请务必参考：\par

            线性空间到域 $V \to \mathbb{F}$ 的函数的集合 $V^*$ ，若满足：\par
            $\forall \phi,\psi \in V^*: (\phi + \psi)(x) = \phi(x) + \psi(x)$\par
            $\forall \phi \in V^*, \forall \alpha \in \mathbb{F}: (\alpha\phi)(x) = \alpha \phi(x)$\par
            构成一个线性空间，称为 $V$ 的\textbf{对偶空间}\index{对偶空间}．对偶空间中的元素是\textbf{线性函数}\index{线性函数}．\par
            > 有限维的情况下，线性空间及其对偶空间维度相等\par
            $\mathbb{R}^n$ 可以视作$n$ 维\textbf{列向量}\index{列向量}表示的实数数组的集合，于是可以用相同维的\textbf{行向量}\index{行向量}表示对偶空间中的\textbf{线性函数}\index{线性函数}．\par
            看一个简单的例子，在域 $\mathbb{F}$ 上，一对2维对偶向量 $x\in V,f\in V^*$ 的共同作用为：\par
            $\ll f,x \gg = f(x) = fx=\begin{bmatrix} f_1 & f_2 \end{bmatrix}\begin{bmatrix} x_1 \\ x_2 \end{bmatrix} =\sum_{i=1}^{2} f^ix_i = f^ix_i$\par
            二元映射 $\ll , \gg$ 是\textbf{双线性映射}\index{双线性映射}（也就是对两个变元都是线性的，易于验证）：\par
            $V^* \times V \to \mathbb{F}$\par
            $f \in V^*$ 作为对偶向量，通过函数的方式作用于 $x \in V$ ．由于是线性函数，总可以写为矩阵乘法的形式．连等式的最后一个表达式去掉了求和符号，这是默认采用了\textbf{Einstein求和规则}\index{Einstein求和规则}，即对上下同时出现的指标项目求和．\par
            注意到，矩阵乘法的形式提示我们，对偶关系是相对的，即 $V^*$ 也是 $V$ 的对偶空间，实际上 $(V^*)^* \sim V$ ，证明略．\par

            我们有了对偶空间和对偶向量的概念后，形象上容易通过列向量和行向量来区分一对对偶向量．而区分它们更重要的因素是向量的坐标如何随着基的变化而变化，也就是\textbf{逆变性}\index{逆变性}和\textbf{协变性}\index{协变性}．\par
            简单起见先来点小学数学．将位移的\textbf{量纲}\index{量纲}记为 $\left[ L \right]$ ．通常我们可以用一个符号 $l$ 指代某个位移，这个位移具有明确的值，它与测量的\textbf{单位}\index{单位}无关，例如：\par
            $l = \left\{ 10 \right\}_{m} = \left\{ 1000 \right\}_{cm}$\par
            基于显然的单位换算：\par
            $1 \text{m} = 100 \text{cm}$\par
            位移的取值与测量单位具有以下的关系：\par
            $l = \left\{ 10 \right\}_{100cm} = \left\{ 1000 \right\}_{1cm}$\par
            稍微啰嗦地展开为：\par
            $l = \left\{ 10 \right\}_{100cm} = 10 \cdot 100cm$\par
            $l = \left\{ 1000 \right\}_{1cm} = 1000 \cdot 1cm$\par
            学过线性代数，我们可以更加名词党地理解为 $1\text{m}$ 和 $1\text{cm}$ 是 $1$ 维位移\textbf{线性空间}\index{线性空间}的两个不同的\textbf{基}\index{基}，而括号里作为测量读数的\textbf{坐标}\index{坐标}是\textbf{逆变}\index{逆变}的．逆变的概念推广到2维：\par
            $\begin{bmatrix} \left\{ 20 \right\}_{100cm} \\ \left\{ 3000 \right\}_{1cm}\end{bmatrix} =\begin{bmatrix} \left\{ 1000 \right\}_{2cm} \\ \left\{ 6 \right\}_{500cm}\end{bmatrix} $\par
            这里提示有两组基，类似前面的也展开：\par
            $\begin{bmatrix} \left\{ 20 \right\}_{100cm} \\ \left\{ 3000 \right\}_{1cm}\end{bmatrix} =\begin{bmatrix} 20 \cdot 100cm \\ 3000 \cdot 1cm \end{bmatrix} =\begin{bmatrix} 100cm & 0 \\ 0 & 1cm \end{bmatrix} \begin{bmatrix} 20 \\ 3000 \end{bmatrix} $\par
            $\begin{bmatrix} \left\{ 1000 \right\}_{2cm} \\ \left\{ 6 \right\}_{500cm}\end{bmatrix} =\begin{bmatrix} 1000 \cdot 2cm \\ 6 \cdot 500cm \end{bmatrix} =\begin{bmatrix} 2cm & 0 \\ 0 & 500cm \end{bmatrix} \begin{bmatrix} 1000 \\ 6 \end{bmatrix} $\par
            类似于前面，当基矩阵变换后，坐标随之变换，变换的方向是相反的．进一步，在 $\mathbb{R}^n$ 中，有基矩阵 $A \to B$ 的变换，满足变换规律：\par
            $B=AT$\par
            那么，坐标向量 $\alpha \to \beta$ 的变换则满足：\par
            $A\alpha = B\beta = AT \beta$\par
            于是\par
            $\alpha = T \beta$\par
            $\beta = T^{-1}\alpha$\par
            这里出现的逆矩阵，就是逆变的含义．\par

            我们又以熟悉的平面旋转为例子．实平面 $\mathbb{R}^2$ 上给定一组基：\par
            $\begin{bmatrix} e_{11} \\ e_{21} \end{bmatrix}, \begin{bmatrix} e_{12} \\ e_{22} \end{bmatrix}$\par
            合写为向量组：\par
            $E=\begin{bmatrix} e_{11} & e_{12} \\ e_{21} & e_{22} \end{bmatrix}$\par
            在 $E$ 下，向量 $v$ 的坐标表示为：\par
            $v=Ex =\begin{bmatrix} e_{11} & e_{12} \\ e_{21} & e_{22} \end{bmatrix} \begin{bmatrix} x_1 \\ x_2 \end{bmatrix}$\par
            取\par
            $\rho(\theta)=\begin{bmatrix} \cos \theta & -\sin \theta \\ \sin \theta & \cos \theta \end{bmatrix} \in SO(2)$\par
            作用于基向量组，得到一组新的基：\par
            $E \cdot \rho(\theta) =\begin{bmatrix} \cos \theta & -\sin \theta \\ \sin \theta & \cos \theta \end{bmatrix} \begin{bmatrix} e_{11} & e_{12} \\ e_{21} & e_{22} \end{bmatrix}$\par
            则向量 $v$ 的坐标表示为：\par
            $v=E \cdot \rho(\theta) \cdot x^\prime =\begin{bmatrix} \cos \theta & -\sin \theta \\ \sin \theta & \cos \theta \end{bmatrix} \begin{bmatrix} e_{11} & e_{12} \\ e_{21} & e_{22} \end{bmatrix} \begin{bmatrix} x_1^\prime \\ x_2^\prime \end{bmatrix}$\par
            基变换前后，向量 $v$ 在自然基下的表达不变：\par
            $Ex = E \cdot \rho(\theta) \cdot x^\prime $\par
            于是\par
            $x^\prime = \rho(\theta)^{-1} \cdot x$\par
            符合逆变的要求．即当基在 $\rho(\theta) \in SO(2)$ 作用下旋转角度 $\theta$ 时，坐标相应反向旋转，符合直观的认知．\par

            回到对偶空间中，一对2维对偶向量 $x\in V,f\in V^*$ 通过双线性映射共同作用：\par
            $\ll f,x \gg = f^ix_i$\par
            给定 $V$ 的一组基 $E = \{e_i\}$ ，则 $V$ 中任意向量可以表为 $v=Ex$ ．这组基 $E$ 可以通过\textbf{Kronecker记号}\index{Kronecker记号}唯一地决定对偶空间 $V^*$ 中的一组基 $E^*=\{e^j\}$ ，满足：\par
            $e^j(e_i)=\delta_{ij}$\par
            其中\par
            $i=j \Rightarrow \delta_{ij}=1$\par
            $i \ne j \Rightarrow \delta_{ij}=0$\par
            这样得到的对偶空间中的基称为\textbf{对偶基}\index{对偶基}．\par
            我们不难验证，由于Kronecker记号的约束，当基以 $T:E \to E^\prime$ 变化时，对偶基相对于基的变化规律是逆变的．对偶向量在对偶空间的自然基（自然基的对偶基）中的表示是不变的，类似于前面讨论的，对偶向量的坐标相对于对偶基也是逆变的．这样，两次逆变的合成，使得对偶向量的坐标相对于 $T$ 的变化是\textbf{协变}\index{协变}的，称为\textbf{协变向量}\index{协变向量}．\par

            对偶空间中，一对对偶向量 $x\in V,f\in V^*$ 通过双线性映射共同作用：\par
            $\ll f,x \gg$\par
            它对两个变元都是线性的．对偶空间可以是\textbf{自对偶}\index{自对偶}的，泛函分析的\textbf{Riesz引理}\index{Riesz引理}给出了构建自对偶空间的方法．自对偶空间的双线性映射是\textbf{正定}\index{正定}且\textbf{共轭对称}\index{共轭对称}的，则称为\textbf{内积}\index{内积}．完备的内积空间称为\textbf{Hilbert空间}\index{Hilbert空间}．在实分析和泛函分析中常见的一个例子是 $L^p$ 空间，即：\par
            设 $\mu$ 是 $\Omega$ 上的测度， $(\Omega,\mathscr{B},\mu)$ 是测度空间，函数 $u(x)$ 满足在 $\Omega$ 上 $\vert u(x) \vert^p$ 可积 $(1\leq p\leq +\infty)$ ，则这类函数的全体称为$(\Omega,\mathscr{B},\mu)$上的 $p$ 次可积函数空间，记作 $L^p(\Omega,\mu)$ ．定义内积：\par
            $\langle u,v \rangle = \int_\Omega u(x) \overline{v(x)} dx$\par
            量子力学中常用的为平方可积的 $L^2$ 空间．\par


\section{自然配对}

作用在$k$-线性空间上，
反变Hom函子相当于线性对偶函子
$(-)^* = h^k = \textbf{Vct}_k(-,k) \in \textbf{Fct}(\textbf{Vct}_k^\text{op},\textbf{Vct}_k)$，
得到$k$-线性泛函的集合，
即$k$-线性对偶空间，
记为$X^* = h^k X = \textbf{Vct}_k(X,k) \in \textbf{Vct}_k$．
再次作用在对偶空间上得到双重对偶空间$X^{**}$．
$x\in X$决定了对偶线性泛函：
$(i_x \in X^{**}): X^* \to k :: \alpha \mapsto i_x(\alpha) = \alpha(x)$，
这样构造了自然嵌入
$X \to X^{**} :: x \to i_x$，
它是$k$-线性单同态：

$$
\begin{aligned}
X^* \times X &\to k \\
(\alpha,x) &\mapsto \alpha x = \alpha(x)
\end{aligned}
\quad
\begin{aligned}
X^{**} \times X^* &\to k \\
(\Phi,\alpha) &\mapsto \Phi \alpha = \Phi(\alpha)
\end{aligned}
\quad
\begin{aligned}
i:X \times X^{**} &\to k \\
x &\mapsto i(x) = x
\end{aligned}
$$

约定$X^* \times X$和$X^{**} \times X^*$二元组的顺序，
以上的双线性形可以记为：
记$\langle \alpha,x\rangle = \alpha(x)$
以及$\langle \Phi,\alpha\rangle = \Phi(\alpha)$，
这样二元组的元就区分为左元和右元．
通常左元称为线性泛函(linear functional)或对偶向量(covector)，
右元仍保留向量(vector)的称谓．


对偶函子的反变性体现在
$f^* = h^k f: Y^* \to X^* :: \beta \mapsto f^*(\beta) = (h^k f)(\beta) = \beta \circ f$，
这样构成了$f\in\textbf{Vct}_k(X,Y)$
的对偶算子$f^*\in\textbf{Vct}_k(Y^*,X^*)$，
进而有双重对偶算子$f^{**}$：

$$
% https://tikzcd.yichuanshen.de/#N4Igdg9gJgpgziAXAbVABwnAlgFyxMJZABgBpiBdUkANwEMAbAVxiRAA0QBfU9TXfIRRkATFVqMWbAJrdeIDNjwEiAFnLj6zVohABrOXyWC1pMdS1TdBnkYEqUIs5sk6Qhhf2VDkTyhdc2D0V7HwA2DQDtNk5bT2MHZAjzCWjdWTiQ7yIADkjUqw4APQAqYK8TFDyUyzdpUvKE8NIARhc090yKxLy2qMLG0KIWsn8Ct3Yi4BKSrkHslBHnfrqpmbmupuGyPvGgrnEYKABzeCJQADMAJwgAWyQyEBwIJBbqBjoAIxgGAAVuoQgK5YY4ACxwIBWbAuHmudwe1GeSCce10F1KAAIADpY744OgYgC8GIAFKCinoMRcAJQknF4ujUonY3EwfEsgDGWCuHKpLKwYAxkxKzPJFyFsJu90QbyeL0QAGYobp6WyCTiBRj6iLieTKbJ3l8fv8trpgWCIXE4dL1HKkABWd4CtxQCBMT4MVjUUEwOhQJBgJgMBiIuhYBhsSBgL2okDonUYvVUyEgD7fP4Atjm8GS+GIFFIxVO6NsV3uz0pn1+yMEGP48M1ktWqVICJ2xAAdmVcZTaeNmbNIJzzbzAE5EfK8rH472jRnTSABdhYLnpUt2y1ZQxnaW3R6Y1X-YhA8HQw3dFHWCO17LCy1bduS7oy-vK76jyeQ08wxGL7XV68Bbyi0SrTmsswAYgbaFjk15IFOd4iHBnYTkgo7IeOG4KgcXBAA
\begin{tikzcd}
X \arrow[dd, "f"'] \arrow[rrrr, "f^* \beta = (h^k f)(\beta) = \beta \circ f \in X^* = h^f X"] &  & {}                                        &  & k \arrow[dd, no head, Rightarrow] &  & X \arrow[dd, "f"'] \arrow[rr] &  & X^* \arrow[rr]                               &  & X^{**} \arrow[dd, "f^{**}"] \\
                                                                                              &  &                                           &  &                                   &  & {} \arrow[rr, Rightarrow]     &  & {} \arrow[rr, Rightarrow]                    &  & {}                          \\
Y \arrow[rrrr, "\beta \in Y^* = h^k Y"']                                                      &  & {} \arrow[uu, "f^* = h^k f"', Rightarrow] &  & k                                 &  & Y \arrow[rr]                  &  & Y^* \arrow[uu, "f^*" description] \arrow[rr] &  & Y^{**}                     
\end{tikzcd}
$$

基于左右元的区分，
一对对偶算子$f$作用在左元，
$f^*$作用在右元．
$\forall x$和$\forall \beta$是独立自由的变量．
若这对对偶算子满足
$ \langle x,f^* \beta \rangle = \langle f x, \beta \rangle $，
则形成了颇为耦合的结构，
称他们相互伴随，
记$f \vdash f^*$．

$$
% https://tikzcd.yichuanshen.de/#N4Igdg9gJgpgziAXAbVABwnAlgFyxMJZABgBpiBdUkANwEMAbAVxiRAA0QBfU9TXfIRRkATFVqMWbdgD0AVN14gM2PASIjy4+s1aIQATUV9VgjaTHUdU-QfnHl-NUJKkAjNsl6QAHR8M6MABzBhgAAgBaUgAzeUiwvwAnQJDWHhMBdRRNDysvNj8A4NCw6MjSCKSU0IcVTJcAFi083TYAD1qnMxQANmaJVv1ojvTHUyzkPssBm18fACMYHDpO8caLT0GQWLkEhaWV0brnIibcme9C6vC2mLi-ReW95OK0pWPuyfdN2avX0rCtz2jzoz2u3HEMCgQXgRFA0USEAAtkgyCAcBAkJoLmxog4EcikABmagYpBuFqzHb4xEoxBNdGYxDEUYEukMsmINys2lIACspKZIh5hMQAsZxJFdL6EsQAHZSXQsAw2Ei6Gg4GSpUgAByCpAATkVytV6s1mO1XLRnJlyxN+jVGq1SjZ5OtTKN6KVKodZud8N5XIpsoVXvtIEd5pporcwc5erDPojfotLsDbnd5ODDCwYG8UAgTHmNWoAAsYHQoGxIHmQMakzW0hQuEA
\begin{tikzcd}
X \arrow[rr, "f"]                             &  & Y                                          &  & x \arrow[rr, maps to]                                                                               &  & fx                                                                 \\
{\langle -,f^* - \rangle} \arrow[u] \arrow[d] &  & {\langle f -,-\rangle} \arrow[u] \arrow[d] &  & {\langle x,f^* \beta \rangle} \arrow[u, maps to] \arrow[d, maps to] \arrow[rr, no head, Rightarrow] &  & {\langle f x, \beta \rangle} \arrow[u, maps to] \arrow[d, maps to] \\
X^*                                           &  & Y^* \arrow[ll, "f^*"]                      &  & f^* \beta                                                                                           &  & \beta \arrow[ll, maps to]                                         
\end{tikzcd}
$$


集合范畴中，
$f\in\textbf{Set}(X,Y)$产生一对函数
$f_*\in\textbf{Set}(\textbf 2^X,\textbf 2^Y):: A\mapsto f_*(A) = f(A) = \text{im} \ f (A)$，
$f^*\in\textbf{Set}(\textbf 2^Y,\textbf 2^X):: B\mapsto f^*(B) = f^{-1}(B) = \text{preim} \ f (B)$．
以幂集上自然的偏序结构为范畴，
反单调Galois连接给出了伴随函子
$f^* \vdash f_*$：

$$
% https://tikzcd.yichuanshen.de/#N4Igdg9gJgpgziAXAbVABwnAlgFyxMJZABgBpiBdUkANwEMAbAVxiRAAUANEAX1PUy58hFACZyVWoxZte-EBmx4CRACwTq9Zq0QcAmnIFLhRAGwap2tgEFDCwcpHIAHBa0zdAMwD6AKgAU1gCUdopCKigAjKSRku46INx8RuFO4rGa0gmhDiYoAMwxcVlsBsn2xhEkpKLFVrpcOZVpNXUeIE2paq2Z9fqdjmY9lu2eAHoBAEIh5WGDKK61ve2TvJIwUADm8ESgngBOEAC2SGQgOBBI4iMJPr4g1Ax0AEYwDOy5ESAMMJ44dgdjqdqBckK4bmwADqQn4ARwe3xebw+zTY+ywmwAFv9yoCTohCudLoh1BCvH4EU9Xu9PiJvr8cfI8UhCaDEJFomSQNC4ZSkTTUbp0VjGXtDviAKwg4kAdmWtz51JRXV0WDA2FgAPFSHMRKQnPiSDATAYDEe-OV83pf0paoSUAgTGePy1QMQurZAE55UaTWbEUraWwfjbHna2A6nS7cdr2Wc2eDDV4JorkUHdCHRSBmXHpVcfboeTB4ebA4KQMLsa78ZFrmyOQXsynS2ny5nq-q68TSUnuTDi6mBSqKxiqzG3VK9YhiOPJXnEM5Z0g5VPREvECv6zOKDwgA
\begin{tikzcd}
PX \arrow[rrrr, "f_*"] \arrow[dd, "\leq"'] &                                                     & {}                                             &                         & PY                                         &  & A \arrow[rr, "f_*"] \arrow[dd, "\leq"'] &  & f_*(A)                                  \\
                                           & X \arrow[rr, "f" description] \arrow[lu] \arrow[ld] & {} \arrow[u, Rightarrow] \arrow[d, Rightarrow] & Y \arrow[ru] \arrow[rd] &                                            &  &                                         &  &                                         \\
PX                                         &                                                     & {}                                             &                         & PY \arrow[llll, "f^*"] \arrow[uu, "\leq"'] &  & f^*(B)                                  &  & B \arrow[ll, "f^*"] \arrow[uu, "\leq"']
\end{tikzcd}
$$

考虑光滑函数$f\in \textbf{Mfd}^\infty(X,Y)$
在$(x,y=f(x))$产生的切空间$(T_x X,T_yY)$
和余切空间$(T^*_x X,T^*_yY)$．
假设$f$诱导了一对$(f_*,f^*)$，
用切向量和余切向量的自然配对来定义${\langle -, - \rangle}$，
构成了推出(pushforward) $f_*$ 和拉回(pullback) $f^*$，
得到类似前述的$f_* \vdash f^*$的伴随结构：

$$
% https://tikzcd.yichuanshen.de/#N4Igdg9gJgpgziAXAbVABwnAlgFyxMJZABgBpiBdUkANwEMAbAVxiRABUB9ADwA0QAvqXSZc+QigBM5KrUYs2XAJ4BNQcJAZseAkQAsM6vWatEtdSO3iiANkNyTbAGacAVAAIaFzaJ0Tk0gCMssYKZgA64Qx0YADmDDDuLh4AtKQp7pEATjHxrEKWYrooBsFG8qYgkdFxCe4A7qROAHqukRAAtjCxdJnhObX5GlpF-nZlDmFVUbl1yZ6k9X0Ded4jfkRkkiEViq08-AU+VsUBpNvljmbs+6prvtYl5ztXIC1t4Z3ddPcnY8+XKbtLo9X6jTakCahSrVWaJNLvdwZbJwwSyGBQWLwIigJxZTpIaQgHAQJAAZkBlWSIGo0QARjAGAAFB7FEAJJw4bx4gmIACs1BJhMpSDATAYDFpdAZzNZEhAWSwsQAFlzBXQsAw2B06Gg4EKjjyOkg7MTSYgKZNTGKJVKZSy-mxFSq1cSNVqzDq9QaNEakAKzUgAJwixA2yXs6WMh3gszO1U0t2a7W6-Wkw3442IQJkQOIENW5ytRP06NytgcrkZ3mmoXZ3PQ0XiiOl2WOuNKhPq5Oe1M+3GZpC5uuBUNvNzcweIAAcgvNAHYx+9J7zAqO88Rq1m13OkPOt0gDHnpwfEEeR6eA3XTQwsGBKlAIEw6QlE8qYHQoGxIPfEzh3d+BD5BQAhAA
\begin{tikzcd}
T_xX \arrow[rr, "f_*"]                        &  & T_yY                                          &  & v \arrow[rr, "f_*"]                                                                                 &  & f_* v                                                           \\
{\langle -,f^* - \rangle} \arrow[u] \arrow[d] &  & {\langle f_* -,- \rangle} \arrow[d] \arrow[u] &  & {\langle w,f^*\omega \rangle} \arrow[u, maps to] \arrow[d, maps to] \arrow[rr, no head, Rightarrow] &  & {\langle f_* v,w \rangle} \arrow[u, maps to] \arrow[d, maps to] \\
T^*_xX                                        &  & T^*_yY \arrow[ll, "f^*"]                      &  & f^*\omega                                                                                           &  & \omega \arrow[ll, "f^*"]                                       
\end{tikzcd}
$$



\section{切空间，余切空间}

平直空间$X = \mathbb R^n$视为仿射空间，
点$x \in X$上所有平移构成线性空间$v \in T_x X$，
函数$f\in C(X)$对方向向量$v\in T_x$的方向导数是：

$$
D_v f \mid_x = \lim_{t \to 0} \frac{f(x + t v) - f(x)}{t} = \text{grad} \ f \cdot v
$$

它是关于$\text{grad} \ f$和$v$的双线性函数，
故$\text{grad} \ f$是$v$的对偶向量，
记$\text{grad} \ f \in T^*_x X$，
方向导数可以写为自然配对$\langle \text{grad} \ f,v \rangle = D_v f \mid_x$．

我们希望把这种$T^*_x X \times T_x X \to \mathbb R$的自然配对结构推广到光滑流形，
在光滑流形的点$x\in X\in\text{Mfd}^\infty$的局部研究微分问题．
令有流形上的光滑函数$f \in C(X)$，
令$\Gamma_x = \{ \gamma: \mathbb R \to X \mid \gamma(0) = x \}$为过$x$的参数曲线集合，
借助参数曲线$\gamma \in \Gamma$，
可以构造一元实可导函数$f\circ\gamma: \mathbb R \to \mathbb R$：

$$
% https://tikzcd.yichuanshen.de/#N4Igdg9gJgpgziAXAbVABwnAlgFyxMJZAJgBoAGAXVJADcBDAGwFcYkQAKADwAIBeHgB1BAc3oBbcfQ7kAlLKGCsYHgA0QAX1LpMufIRTlSxanSat2HHPx5zFyxVJwALAEaueAJU3aQGbHgERAAsxqYMLGyIIABm3ArCDsJObh7eGqYwUCLwRKAxAE4Q4kgAjDQ4EEhGZpHswmKS9CAV9FiM7FJocJU++UUliDW9iGS1FtExLSA4bR3RXT1VWv3FZRVVozQRE7GKAMZYBfuKjVLTjPSuMIwACrqBBiAFWCLOONOz7Z303b0ZGiAA
\begin{tikzcd}
                                                                                               &  & (x = \gamma(0)) \in X \arrow[rrdd, "f", maps to] &  &                    \\
                                                                                               &  &                                                  &  &                    \\
(t = 0) \in \mathbb R \arrow[rruu, "\gamma", maps to] \arrow[rrrr, "f \circ \gamma"', maps to] &  &                                                  &  & f(x) \in \mathbb R
\end{tikzcd}
$$

参数曲线$\gamma$和光滑函数$f\in C(X)$在$x$点附近的局部行为，
而光滑函数$f\in C(X)$在$x$上有茎$[f] \in \mathcal O_x$，
两者在$x$附近的局部作用描述为双线性映射：

$$
\begin{aligned}
  \mathcal O_x \times \Gamma_x &\to \mathbb R \\
  ([f],\gamma) &\mapsto \langle [f],\gamma \rangle = \frac{\text d (f \circ \gamma)}{\text d t}
\end{aligned}
$$


\section{局部线性化}
            本讲回顾了高等数学/数学分析中经典的Taylor展开技术．我们关注的是Taylor展开的线性部分，这部分的线性系数就是函数的导数．在多元函数中，这部分的一组线性系数构成了某种线性组合．\par
            站在线性代数的角度看待Taylor展开的线性部分，我们考虑两个问题：\par
            坐标的微分能否构成某种线性空间的向量？\par
            坐标的微分能否构成这个线性空间的基？\par
            我们研究了坐标分量函数，这类函数的微分就是坐标的微分，几何上它们构成了一组正交的超平面系．借助这类函数，我们发展出了微分形式的概念，确认了以上两个问题，即坐标的微分就是对偶空间的向量，坐标的微分正是这个线性对偶空间的基．\par
            之后我们研究了这类微分形式的几何意义，通过它们线性组合，可以构成具有任意梯度的微分形式．\par
            上一讲中，我们在考虑如何把 $\mathbb{R}^n$ 上的向量场推广到一般的光滑流形 $M$ 时，遇到了一些问题．这些问题反映出流形和平直空间内在的不同．这些问题的解决需要等到完全理解\textbf{切空间/余切空间}\index{切空间/余切空间}的概念．在展开讲解之前，在这一讲中我们先谈谈局部线性化的思想，同时储备一些技术工具．\par
            回顾微分流形 $M$ 的定义，其中提到\par
            > $\forall p \in M$ 都 $\exists U$ ： $p\in U \subset M$ 使得开邻域 $U$ 和 $\mathbb{R}^n$ 中的开集同胚．\par
            这使我们注意到：\par
            1. 这里是\textbf{逐点定义}\index{逐点定义}的，且只考虑每一点的某个\textbf{邻域}\index{邻域}的存在性．这提示我们采用了\textbf{局部}\index{局部}的方法．\par
            2. 局部（开邻域）和 $\mathbb{R}^n$ 中的开集同胚，意味着在局部试图接近 $\mathbb{R}^n$ 这样的平直空间，所谓平直，就是\textbf{线性}\index{线性}，实际上 $\mathbb{R}^n$ 中嵌入的平直空间（直线、平面、超平面）就称为\textbf{线性流形}\index{线性流形}．\par
            这些思路，我们总结为\textbf{局部线性化}\index{局部线性化}．局部线性化的思想贯穿于微分几何的研究中，这也是我们一直强调同学们学好线性代数的原因（另一个思想是借助拓扑学的整体微分几何）．这一讲中，我们详细谈一下如何借助古老的Taylor展开技术，把局部线性化的方法发展为微分形式．\par

            微积分中我们熟知\textbf{Taylor展开}\index{Taylor展开}，在一维的情况，函数 $f(x)$ 在点 $x_0 \in \mathbb{R}$ 附近可以写为幂级数和的形式：\par
            $f(x) = \sum_{k=0}^\infty { \frac{f^{(k)}(x_0)}{k!}(x-x_0)^k } $\par
            $= f(x_0) + \frac{f^\prime(x_0)}{1!}(x-x_0) + \frac{f^{\prime\prime}(x_0)}{2!}(x-x_0)^2 + \dots$\par
            注意到 $k=0$ 次项 $f(x_0)$ 是常数项，而当 $x \to x_0$ 求极限时：\par
            $\frac{df}{dx}(x) \Bigg\vert_{x=x_0}   =\lim_{x \rightarrow x_0} { \frac{f(x)-f(x_0)}{x-x_0} }$\par
            $=\lim_{x \rightarrow x_0} \frac{1}{x-x_0} \Bigg\{  \frac{f^\prime(x_0)}{1!}(x-x_0) + \frac{f^{\prime\prime}(x_0)}{2!}(x-x_0)^2 + \dots \Bigg\}$\par
            在 $x_0$ 处求导将消去 $1$ 次以上高阶无穷小量，只剩下 $1$ 次的线性项，得到我们熟知的关系：\par
            $\frac{df}{dx}(x) \Bigg\vert_{x=x_0} = f^\prime(x_0)$\par
            站在我们现在的角度看，Taylor展开就是一种函数逼近技术，当考察固定点 $x_0 \in \mathbb{R}$ 附近的局部（邻域）时，成为了一种\textbf{线性化}\index{线性化}的，只保留一次导数 $f^\prime(x_0)$ 的逼近．\par
            注意到Taylor展开的局部线性化也可以表示成：\par
            $df=f^\prime dx$\par
            这一方面反映了导数的局部线性的性质，另一方面，导数 $f^\prime$ 在固定点上是常数，它本身可以视为\textbf{线性系数}\index{线性系数}．在 $n$ 维看得更加清楚．$n$ 维函数 $f(x)$ 在点 $x_0 \in \mathbb{R}^n$ 附近的Taylor展开，进行局部线性化后得到：\par
            $df=\sum_{k=1}^n \frac{\partial f}{\partial x_i} dx^i =\frac{\partial f}{\partial x_i} dx^i =(\partial_i f) dx^i$\par
            这里用了Einstein求和约定．显然，偏导数 $\{ \partial_i f \}$ 起到了\textbf{线性系数}\index{线性系数}的作用．现在我们得到了某种类似\textbf{线性组合}\index{线性组合}的关系．然而，若要放心地用线性代数的方法来处理，我们首先需要明确两个问题：\par
            1. $df$ 和 $\{dx^i\}$ 是否是某个向量空间中的向量？\par
            2. 如果是， $\{dx^i\}$ 能否构成基？\par
            如果我们能够解决这两个问题，过去我们所理解的所谓\textbf{无穷小量}\index{无穷小量}的\textbf{微分}\index{微分}将被赋予线性空间的新的内涵，它是将来要讨论的\textbf{余切向量}\index{余切向量}的原形．\par

            在 $\mathbb{R}^n$ 中，记自然基为 $\{x_i\}$ ．考虑一组特殊的函数\textbf{分量函数}\index{分量函数} $C=\{x^i(x)=x_i\}$ ，其中每一个表现为：\par
            $x^i:\mathbb{R}^n \to \mathbb{R}$\par
            $x \mapsto x^i(x) = x_i$\par
            即把点 $\forall x\in\mathbb{R}^n$ 的第 $i$ 个分量 $x_i$ 输出为函数值．\par
            > 注意：指标 $i$ 在上面表示它是\textbf{函数}\index{函数}： $x^i=x^i(x)$  指标在下面表示它是坐标\textbf{分量}\index{分量}： $x_i \in \mathbb{R}$\par
            分量函数 $x^i=x^i(x)$ 对分量 $x_j$ 求偏导数：\par
            $\frac{dx^i(x)}{dx_j}=\delta_{ij}$\par
            这里用了Kronecker记号．于是仅在指标 $i=j$ 时有意义：\par
            $\frac{dx^i(x)}{dx_i}=\delta_{ii} = 1$\par
            即\par
            $dx^i=dx^i(x)=dx_i$\par
            上式中，左边的指标 $i$ 在上面表示函数的微分，右边的指标在下面表示自变量的微分．我们理解为，\textbf{函数微分}\index{函数微分}和\textbf{自变量微分}\index{自变量微分}可以等同看待．\par
            由分量函数的函数微分和自变量微分的等同关系可见，分量函数 $C=\{x^i(x)=x_i\}$ 的几何意义是，每一个函数 $x^i(x)=x_i$ 都是线性函数，它构成一个\textbf{超平面（线性流形）}\index{超平面（线性流形）}，这组超平面相互\textbf{正交}\index{正交}．\par

            将 $n$ 维实线性函数所构成的空间记为 $L(\mathbb{R}^n)$ ．这个线性函数空间蕴含了自然的线性运算，即 $\forall \alpha,\beta \in \mathbb{R}, \forall f, g \in L(\mathbb{R}^n)$ ：\par
            $f, g: \mathbb{R}^n \to \mathbb{R}$\par
            $(\alpha f + \beta g)(x) = \alpha f(x) + \beta f(x), \forall x \in \mathbb{R}^n$\par
            我们注意到上面讨论的分量函数\par
            $C=\{x^i(x)=x_i\}$\par
            中的每个函数都是 $n$ 维实线性函数．前面谈到它们构成了正交的超平面，代数上意味着 $C$ 是函数空间 $L(\mathbb{R}^n)$ 中的\textbf{线性无关}\index{线性无关}子集．\par
            现在考虑某一点 $x_0 \in \mathbb{R}^n$ 附近的线性函数空间 $L_{x_0}(\mathbb{R}^n)$ ．如前面的讨论，任意函数 $f(x)$ 在点 $x_0$ 附近若可微，则可以通过Taylor展开，在 $x_0$ 附近消去高阶无穷小量后，局部视为一个线性函数 $f \in L_{x_0}(\mathbb{R}^n)$ ：\par
            $f: \mathbb{R}^n \to \mathbb{R}$\par
            $x \mapsto f(x)=f(x_0) + \lambda_i x^i$\par
            更重要的，可以把这个线性函数 $f$ 表达为点 $x_0$ 附近分量函数\par
            $C=\{x^i(x)=x_i\} \subset L_{x_0}(\mathbb{R}^n)$\par
            的线性组合：\par
            $f(x)=f(x_0) + \lambda_i x^i =f(x_0) + \lambda_i x^i(x) $\par
            对它求全微分：\par
            $df=\sum_{k=1}^n \frac{\partial f}{\partial x_i} dx^i =\frac{\partial f}{\partial x_i} dx^i =(\partial_i f) dx^i = \lambda_i dx^i$\par
            这里的 $\lambda_i  = \partial_i f$ 是线性系数．这里出现的 $dx^i$ 既可以视为变量 $x^i$ 的微分，也可以视为函数 $x^i(x)=x_i$ 的微分．\textbf{非常重要！}\index{非常重要！}回到前面我们提出的问题：\par
            1. $df$ 和 $\{dx^i\}$ 是否是某个向量空间中的向量？\par
            2. 如果是， $\{dx^i\}$ 能否构成基？\par
            现在，在某一点 $x_0 \in \mathbb{R}^n$ 附近，由于 $f=f(x)$ 被局部线性化， $f$ 和 $\{x^i(x)=x_i\}$ 的成员都是 $L_{x_0}(\mathbb{R}^n)$ 中的元素．我们把微分算子 $d$ 视为从\textbf{线性函数空间}\index{线性函数空间}到\textbf{函数微分集合}\index{函数微分集合}的映射：\par
            $d:L_{x_0}(\mathbb{R}^n) \to D_{x_0}(\mathbb{R}^n)$\par
            $f \mapsto df$\par
            $x^i = x^i(x) \mapsto dx^i$\par
            由于微分的\textbf{线性性}\index{线性性}，自然可以在函数微分集合 $D_{x_0}(\mathbb{R}^n)$ 中构建线性运算，构成线性的\textbf{函数微分空间}\index{函数微分空间}．于是，$df$ 和 $\{dx^i\}$ 是函数微分线性空间 $D_{x_0}(\mathbb{R}^n)$ 中的向量．\par
            其次，由于 $C=\{x^i(x)=x_i\} \subset L_{x_0}(\mathbb{R}^n)$ 是一个线性无关的集合，可以证明微分算子 $d$ 将它映射成函数微分线性空间 $D_{x_0}(\mathbb{R}^n)$ 中的基$\{dx^i\}$．\par
            最后，根据全微分公式：\par
            $df=(\partial_i f) dx^i = \lambda_i dx^i$\par
            即：任意函数 $f(x)$ 在点 $x_0 \in \mathbb{R}^n$ 附近若可微，则可以通过Taylor展开在 $x_0$ 附近进行局部线性化，得到函数微分 $df=(\partial_i f) dx^i$ ，它是基$\{dx^i\}$ 的线性组合．\par
            这个函数微分线性空间 $D_{x_0}(\mathbb{R}^n)$ 中的元素称为\textbf{1-微分形式(differential form)}\index{differential form}，简称\textbf{1-形式}\index{1-形式}．本质上，1-形式相当于局部（在相差一个常数的意义下等价，后面详述）线性函数．\par

            函数 $C=\{x^i(x)=x_i\}$ 的几何意义是，在点 $x_0 \in \mathbb{R}^n$ 附近，每一个函数 $x^i(x)=x_i$ 都是 $L_{x_0}(\mathbb{R}^n)$ 中的线性函数，它构成一个\textbf{超平面（线性流形）}\index{超平面（线性流形）}，这组超平面相互\textbf{正交}\index{正交}．\par
            微分算子 $d$ 将这组线性函数映射成函数微分，即：\par
            $d:L_{x_0}(\mathbb{R}^n) \to D_{x_0}(\mathbb{R}^n)$\par
            $x^i = x^i(x) \mapsto dx^i$\par
            注意这种映射不是单射，若有任意实数 $c\in \mathbb{R}$ 则：\par
            $x^i(x) \mapsto dx^i$\par
            $x^i(x) + c \mapsto d(x^i + c) = dx^i$\par
            即，相差为常数的线性函数具有相同的函数微分，相同的函数微分对应着等价类．由于任意线性函数也是 $C=\{x^i(x)=x_i\}$ 的线性组合，这种等价关系对于任意函数微分$df \in D_{x_0}(\mathbb{R}^n)$ 都成立．于是满足微分为 $df$ 的线性函数是一个集合，这个集合中的线性函数两两之间相差一个常数．几何上看， $df$ 决定了一组平行的（相差常数）超平面（线性函数）．\par
            作为1-形式， $df=(\partial_i f) dx^i$ 体现了基$\{dx^i\}$ 的线性组合，几何上看，每一个基向量（1-微分形式）都代表一组平行超平面，基向量之间相互正交，通过系数（偏导数 $\{\partial_i f\}$ ）组合，构成了具有任意\textbf{梯度}\index{梯度}的一组平行超平面．\par
            \par


\section{微分形式}

余切向量$\text d f$使用函数微分的记号是有意为之．
前面讨论了$\text d f = \text{grad} \ f$相当于是流形上的梯度．
梯度作为线性算子

$$
\begin{aligned}
\text d: C(X) &\to T^*X \\
f &\mapsto \text d f  
\end{aligned}
$$

线性空间$V\in\textbf{Vct}$可以构造$i$-阶张量积$T^iV = V \otimes \cdots \otimes V$，
以及\textbf{张量代数}(tensor algebra) $TV = \oplus_{i=0}^\infty T^i V$，
向量以嵌入方成为张量代数的元$\iota: V \to TV$．

为推广$\text d x \text d y = \text d y \text d x$所体现的对称性，
考虑$TV$中如下生成的理想

$$
I = \langle \{ u \otimes v - v \otimes u \mid u,v \in V \} \rangle
$$

有到商代数的自然投射$\pi: TV \to SV = TV/I$，
从而有复合同态$\varphi = \pi\circ\iota$：

$$
\begin{aligned}
  \varphi: V &\to SV = TV/I \\
  u &\mapsto \varphi(u) = u + I \\
  v &\mapsto \varphi(v) v + I \\
\end{aligned}
$$

计算$\varphi(u) \otimes \varphi(v) = u \otimes v + I,\ \varphi(v) \otimes \varphi(u) = v \otimes u + I$，
注意到$u \otimes v - v \otimes u \in I$，
求差$\varphi(u) \otimes \varphi(v) - \varphi(v) \otimes \varphi(u) = (u \otimes v - v \otimes u) + I = 0$，
于是商代数$SV$具有对称性，
成为$V$的\textbf{对称代数}(symmetric algebra)．

为推广$\text dx \text dx = 0$所体现的“外”性质，
考虑$TV$中如下生成的理想

$$
I = \langle \{ v \otimes v \mid v \in V \} \rangle
$$

有到商代数的自然投射$\pi: TV \to \wedge V = TV/I$．
张量代数$(TV,\otimes)$上的张量积乘法，
在自然投射下成为$(\wedge V = TV/I,\wedge)$上的外积：

$$
u \wedge v = u \otimes v + I
$$

计算$v \wedge v = v \otimes v + I$，
注意到$v \otimes v \in I$，
故$v \wedge v = 0$，
于是商代数$\wedge V = TV/I$具有外性质，
成为$V$的\textbf{外代数}(exterior algebra)．

在光滑流形$X$上，
微分形式代数$(\Omega^\bullet(X),\wedge,\text d)$就是这样的外代数．


\section{Riemann流形}

            现在我们要开始接着过去的系列\par
            - 几何与物理I：流形上的分析力学\par
            - 几何与物理II：电动力学\par
            继续同步学习几何和物理．在过去两个系列中我们掌握了大量微分流形、拓扑和张量方面的知识，这些知识是继续本系列 \textbf{几何与物理III：Riemann几何与广义相对论}\index{几何与物理III：Riemann几何与广义相对论} 的基础．众所周知科学史上对广义相对论的研究伴随着Riemann几何中宝藏的挖掘，而Riemann几何的数学结构是附加在光滑流形上的．我们在这一系列中将简要地给出Riemann几何相关的重要概念，并试图让读者对广义相对论有一定的物理直观．这一系列完成后，我们在将来讨论相对论量子力学、量子场论等有关统一理论的主题时，读者将不至于缺乏最基本的概念．\par
            我们在过去分析力学系列、电动力学系列中都看到了相当程度的几何化．MP系列的特点就是尽可能多用先进的几何知识描述老的物理问题．在量子力学的系列，我们学到的是完全不一样的数学工具——主要是泛函分析和其中的算子理论．然而，注意到在量子力学中我们讨论的是Hilbert空间上的算子，而Hilbert空间建立的基础是内积．内积就是抽象空间中的几何化工具，有了内积才能讨论投影、正交、直和、分解等概念．所以，即使是在量子力学系列中，我们仍然隐含着几何化物理问题的主脉．\par
            现在，我们来到了\textbf{广义相对论(general relativity)}\index{general relativity}它是相当几何化的一个物理分支．在广义相对论中，因为质量（引力、动量等）弯曲的时空流形被赋予度量/度规，有了度规就可以确定长度．我们将特别关注切丛、联络——在弯曲时空流形上的平行移动．在切丛上有一种叫做Levi-Civita的联络，可以保持度规且是无挠的．实际上任何的度规都可以确定这样的联络，它成为Riemann流形的切丛上最合适的联络，而这个联络的曲率就是Riemann曲率张量．\par
            对Riemann曲率张量进行处理后，可以得到Ricci张量、广义相对论的Lagrange量，以及最终的Einstein张量，这是Einstein广义相对论方程的一半．Einstein方程描述的是时空受到物质（或者任何有能量或动量的事物）而弯曲的程度．类似于过去我们用外微分描述Maxwell方程的方式，我们要用Bianchi恒等式和能动张量(stress-energy tensor)来构造Einstein的广义相对论方程，这个方程最基本的解就是球对称静态真空的Schwarzschild解．\par
            本讲先介绍Riemann几何最基本的要素——Riemann度量，或者叫做度量张量，在物理上叫做度规张量．\par

            过去讲到Minkowski空间的时候，介绍过Riemann度量\par
            [MP35：张量专题(4)：Riemann度量、Minkowski空间、Lorentz张量](https://zhuanlan.zhihu.com/p/46700297)\par
            微分流形 $M$ 上的 $(0,2)$ 型张量 $g$ 若在 $\forall p \in M$ 上，对于切向量 $\forall u, v \in T_pM$ 都满足：\par
            $g_p(u,v) = g_p(v,u)$\par
            $g_p(u, u) \geq 0$\par
            $g_p(u, u) = 0 \Leftrightarrow u=0$\par
            则它是个对称正定双线性型，称为\textbf{Riemann度量(Riemannian metric)}\index{Riemannian metric}，又叫\textbf{度量张量(metric tensor)}\index{metric tensor}，物理上叫\textbf{度规张量}\index{度规张量}．\par
            这种 $(0,2)$ 型张量容易让我们联想到内积，内积是一对对偶空间上的双线性映射\par
            $\langle \cdot, \cdot \rangle: T_p^*M \times T_pM \to \mathbb R$\par
            而通过Riesz引理可以过渡到自对偶空间上的双线性映射\par
            $\langle \cdot, \cdot \rangle: T_pM \times T_pM \to \mathbb R$\par
            而 $(0,2)$ 型张量(场) $g$ 本质上可以写成如下的逐点映射：\par
            $g_p: T_pM \otimes T_pM \to \mathbb R$\par
            选定一个局部坐标系 $\{x^k\}$ 后，Riemann度量可以用对偶基的张量积来表示\par
            $g_p = g_\mu^\nu(p) dx^\mu \otimes dx^\nu$\par
            微分流形 $M$ 上的张量场可以写为\par
            $g_\mu^\nu(p) = g(\partial_\mu,\partial_\nu)$\par
            它可以用矩阵表示为\par
            $[g_\mu^\nu] \in \mathbb R^{n \times n}, \ \ n =\text{dim}(M)$\par
            对于无穷小的位移有\par
            $ds^2 = g(dx^\mu\partial_\mu, dx^\nu\partial_\nu) = dx^\mu dx^\nu g(\partial_\mu, \partial_\nu) = g_\mu^\nu dx^\mu dx^\nu$\par
            这种正定二次形也相当于是度量．\par

            我们也可以从古典微分几何来理解Riemann度量．最简单最平凡的情况还是欧式空间的自然坐标，它可以通过简单的Cartesian积构造：\par
            $\mathbb{R}^n = \mathbb{R} \times \mathbb{R} \times \dots \times \mathbb{R}$\par
            所以自然坐标系是我们永远要特别关注的坐标系．以三维空间 $\mathbb{R}^3$ 为例，有自然坐标系$(x^1,x^2,x^3)$，常常为了方便引入另外一个\textbf{曲线坐标系}\index{曲线坐标系}，譬如球坐标系 $(r,\theta,\varphi)$ ，用曲线坐标系 $(u^1,u^2,u^3)$ 表示：\par
            $\begin{cases} u^1 = r \\ u^2 = \theta \\ u^3 = \varphi \end{cases}$\par
            其Jacobi行列式不为零\par
            $\text{det}\Big[\frac{\partial u^i}{\partial x_i}\Big] \neq 0 $\par
            故两个坐标系的点以可微的方式对应．\par
            点作为向量可以表示为\par
            $\textbf{r} = \textbf{r}(x^1,x^2,x^3) = \textbf{r}(u^1,u^2,u^3)$\par
            在自然坐标下，我们常常考虑\textbf{偏导数}\index{偏导数}，对应地在曲线坐标下我们考虑\par
            $\textbf{r}_i = \frac{\partial\textbf{r}}{\partial u^i}$\par
            这是坐标曲线在点上的\textbf{切向量}\index{切向量}，这组切向量构成曲线坐标的\textbf{自然标架}\index{自然标架}．\par

            自然基中的点 $(x^1,x^2,x^3)$ 视为对于原点具有长度 $r$ ，于是在自然标架上的投影分量为：\par
            $\begin{cases} r_1 = \sin\theta\cos\varphi\textbf{i} + \sin\theta\sin\varphi\textbf{j} + \cos\theta\textbf{k} \\ r_2 = r\cos\theta\cos\varphi\textbf{i} + r\cos\theta\sin\varphi\textbf{j} - r\sin\theta\textbf{k} \\ r_3 = -r\sin\theta\sin\varphi\textbf{i} + r\sin\theta\cos\varphi\textbf{j} \end{cases}$\par
            现在考虑一个球坐标下的变化量微分\par
            $(dr,d\theta,d\varphi)$\par
            它的长度需要用自然坐标系计算，且满足\par
            $ds^2 = dr_1^2 + dr_2^2 + dr_3^2$\par
            根据上式可以得到\par
            $ds^2 = dr^2 + r^2d\theta^2 + r^2\sin^2\theta d\varphi^2$\par
            如果读者拥有良好的数学感觉，可以发现两个坐标系下都具有正定二次型的形式．既然我们是来研究张量的，那必然要用线性代数的方式来思维，于是给出标准的表示方法：\par
            $ds^2 =  \begin{bmatrix} dr_1 & dr_2 & dr_3 \end{bmatrix} \begin{bmatrix} 1 & & \\ & 1 & \\ & & 1 \end{bmatrix} \begin{bmatrix} dr_1 \\ dr_2 \\ dr_3 \end{bmatrix}$\par
            $=  \begin{bmatrix} dr & d\theta & d\varphi \end{bmatrix} \begin{bmatrix} 1 & & \\ & r^2 & \\ & & \sin^2\theta \end{bmatrix} \begin{bmatrix} dr \\ d\theta \\ d\varphi \end{bmatrix}$\par
            这个变换矩阵\par
            $[g_i^j] = \begin{bmatrix} 1 & & \\ & r^2 & \\ & & \sin^2\theta \end{bmatrix}$\par
            就是曲线坐标下的度量．\par

            现在我们探讨一下在非自然坐标系下对向量的微分．自然基展开的向量 $\textbf{v} = v^k\textbf{r}_k$ 进行微分，按照Leibniz法则：\par
            $d\textbf{v} = d(v^k\textbf{r}_k) = (dv^k)\textbf{r}_k + v^kd\textbf{r}_k$\par
            RHS的第一项 $(dv^k)\textbf{r}_k$ 是自然坐标的微分，第二项 $v^kd\textbf{r}_k$ 是一个补偿项，它补偿了自然标架本身的变化．\par
            $v^kd\textbf{r}_k = v^k\big( (\partial_j\textbf{r}_k) du^j \big)$\par
            上讲提到用自然坐标展开有：\par
            $\partial_j\textbf{r}_k = \Big\{\begin{matrix} i \\ jk \end{matrix}\Big\} \textbf{r}_i$\par
            于是\par
            $v^kd\textbf{r}_k = v^k\big( \Big\{\begin{matrix} i \\ jk \end{matrix}\Big\} \textbf{r}_i du^j \big) = \Big(\Big\{\begin{matrix} i \\ jk \end{matrix}\Big\} v^k du^j \Big) \textbf{r}_i$\par
            那么整个微分可以写为自然坐标的展开：\par
            $d\textbf{v} = (dv^i)\textbf{r}_i + v^kd\textbf{r}_k  = (dv^i)\textbf{r}_i + \big( \Big\{\begin{matrix} i \\ jk \end{matrix}\Big\} v^k du^j \big) \textbf{r}_i$\par
            $= \Bigg( dv^i + \Big\{\begin{matrix} i \\ jk \end{matrix}\Big\} v^k du^j \Bigg) \textbf{r}_i $\par
            其中的系数项就是普通的微分和协变的补偿项．我们进一步还可以考虑普通微分的系数对坐标的展开：\par
            $dv^i + \Big\{\begin{matrix} i \\ jk \end{matrix}\Big\} v^k du^j  = \partial_j v^i du^j + \Big\{\begin{matrix} i \\ jk \end{matrix}\Big\} v^k du^j $\par
            $= \Big( \partial_j v^i + \Big\{\begin{matrix} i \\ jk \end{matrix}\Big\} v^k \Big) du^j $\par
            这样就变成：\par
            $d\textbf{v} = \Bigg( dv^i + \Big\{\begin{matrix} i \\ jk \end{matrix}\Big\} v^k du^j \Bigg) \textbf{r}_i $\par
            $= \Bigg( \Big( \partial_j v^i + \Big\{\begin{matrix} i \\ jk \end{matrix}\Big\} v^k \Big) du^j \Bigg) \textbf{r}_i $\par
            从而可以构造一个微商也就是导数：\par
            $D_j v^i = \partial_j v^i + \Big\{\begin{matrix} i \\ jk \end{matrix}\Big\} v^k$\par
            称为\textbf{协变导数(covariant derivative)}\index{covariant derivative}．\par
            上讲谈到，在自然坐标系中，度量$g_i^j$是常数不变的，故Christoffel记号为零，于是有以下的对易性：\par
            $D_iD_j - D_jD_i = 0$\par
            这使我们联想到量子力学中的对易子，对易子就是通过交换顺序后求差来衡量不对易程度的，和这里如出一辙．\par

            很久以来我们都强调一个观点，就是向量生而是为函数求方向导数的：\par
            [MP9：重新认识向量场：方向导数、偏导数、梯度](https://zhuanlan.zhihu.com/p/44115400)\par
            因此，在微分流形 $M$ 上的向量场 $X$ 若是和光滑函数 $f \in C^\infty(M)$ 写在一起，需要把向量场理解为方向导数算子——它是一个我们在泛函分析和量子力学中长期研究的线性算子：\par
            $X: C^\infty(M) \to C^\infty(M)$\par
            $f \mapsto Xf = X(f) = D_X f$\par
            现在考虑Riemann流形上的向量场．两个向量场 $U,V$ 通过Riemann度量张量 $g$ 的作用，产生一个标量场：\par
            $g: TM \times TM \to C^\infty(M)$\par
            $(U,V) \mapsto g(U,V)$\par
            而第三个向量场作为方向导数算子的作用：\par
            $W: C^\infty(M) \to C^\infty(M)$\par
            $g(U,V) \mapsto Wg(U,V) = D_W g(U,V)$\par
            相当于\par
            $TM \times (TM \times TM) \to C^\infty(M)$\par
            $\big(W, (U,V)\big) \mapsto W g(U,V) = D_W g(U,V)$\par
            用大白话理解：两个向量场$U,V$ 在Riemann度量张量作用下产生的标量场是一个逐点定义的函数，这个函数被向量场 $W$ 逐点作用．\par
            我们记得前面说过Levi-Civita是由Riemann度量 $g$ 确定的，在切丛 $TM$ 上唯一一个保持度量且无挠的联络．令 $\nabla$ 为Riemann流形上的Levi-Civita联络，它可以把向量场$U,V$ 按照向量场 $W$ 的方向得到向量场 $\nabla_W U, \nabla_W V$ ，而且能保持向量的度量，即：\par
            $Wg(U,V) = g(\nabla_W U, V) + g(U, \nabla_W V)$\par
            量子力学有算符的对易子，微分几何中也有向量场的Lie括号．微分流形 $M$ 上的向量场 $U,V$ 可以对光滑函数 $f \in C^\infty(M)$ 求复合方向导数，相当于是一种复合对易的算子，通过\textbf{Lie括号(Lie bracket)}\index{Lie bracket}给出：\par
            $[U,V$](f) = [UV - VU](f) = U(V(f)) - V(U(f))\par
            $= D_U(D_Vf) - D_V(D_Uf) = (D_UD_V - D_VD_U)f$\par
            它就是对易子．Levi-Civita联络还可以确保在平行移动过程中的\textbf{无挠(torsion free)}\index{torsion free}，即：\par
            $[U,V] = \nabla_U V - \nabla_V U$\par

            现在开始真正的广义相对论．在狭义相对论中，我们了解了相对惯性（匀速）运动的参考系之间的Lorentz变换和描述它的Minkowski空间．现在，我们要考虑更进一步推广的Riemann流形和其上的Levi-Civita联络．\par
            [MP60：平行移动与Levi-Civita联络](https://zhuanlan.zhihu.com/p/46700218)\par
            [MP61：Riemann流形上的向量场](https://zhuanlan.zhihu.com/p/46720067)\par
            前面已经讨论了许多Levi-Civita联络的性质，我们知道在Levi-Civita联络上进行的向量场的平行移动是保持度量且无挠的．进一步，我们把Levi-Civita联络上的曲率定义为\textbf{Riemann曲率张量(Riemann curvature tensor)}\index{Riemann curvature tensor}．\par
            本文首先给出Lie括号/对易子在平直空间的向量场上的作用，了解到Lie括号计算的是复合方向导数的不对易（不交换）程度．在平直空间中的偏导数之间是对易的，而对于任意的向量场，这一点不是显然的．\par
            进一步，我们在弯曲的空间上考虑联络的曲率，特别是Riemnn流形上的Levi-Civita联络上的曲率．我们注意到，曲率的前两项相当于是某种形式的对易子，而第三项是一种修正．\par
            Riemann曲率张量是(1,3)型张量，其缩并成为(0,2)型Ricci张量，进一步可以构造出Einstein张量，它是广义相对论的Einstein场方程的重要构成部分．\par

             微分流形 $M$ 上的向量场 $U,V$可以对光滑函数 $f \in C^\infty(M)$ 求复合方向导数，通过\textbf{Lie括号(Lie bracket)}\index{Lie bracket}给出：\par
            $[U,V$](f) = [UV - VU](f) = U(V(f)) - V(U(f))\par
            $= D_U(D_Vf) - D_V(D_Uf) = (D_UD_V - D_VD_U)f$\par
            作为方向导数算符的运算，它就是对易子：\par
            $[U,V$] = UV - VU\par
            Lie括号计算的是复合方向导数的不对易（不交换）程度．在平直空间中的偏导数之间是对易的，即：\par
            $[\partial_i, \partial_j$] = 0\par
            这是我们在量子力学中所熟知的．Lie括号展开后可以理解为两次复合导数，所谓对易即是\par
            $\partial_i\partial_j = \partial_j\partial_i$\par
            按照两个不同偏导数的方向，最后可以在右上角相交，成为交换的\par
            这在多元微积分中是显然的，所以我们在学习多元微积分的时候从未考虑过复合偏导数的顺序问题．然而，对于任意的向量场，这一点不是显然的：\par
            不对易，右上角没有相交，而是产生了一个偏差\par
            上图的右上角所产生的偏差就是Lie括号．\par

            弯曲的空间具有曲率，而曲率是一种内在的属性，不必借助包含它的空间来观察．简单地说，如果让向量从一点出发，按照平行移动的方式经过一个路径后回到这一点，如果空间是弯曲的，那么向量回到这一点后可能发生旋转．\par
            沿着A-N-B-A进行平行移动，回到起点后产生了夹角(wiki)\par
            在流形 $M$ 的闭环路 $\gamma$ 上进行平行移动，回到起点 $p \in M$ ，相当于一个变换：\par
            $H(\gamma, D): TpM \to TpM$\par
            其中 $H(\gamma, D)$ 是由路径 $\gamma$ 和联络 $D$ 决定的，称为\textbf{和乐(holonomy)}\index{holonomy}．用代数的方法研究和乐，通过类似于基本群的方式构造路径的自由群，形成\textbf{和乐群(holonomy group)}\index{holonomy group}．\par

            Riemann流形上有向量场 $U,V,W \in TM$ ，在Levi-Civita联络 $\nabla$ 的作用下，可以产生\par
            $\nabla: TM \times TM \to TM$\par
            $(U,V) \mapsto \nabla_U V$\par
            $(V,U) \mapsto \nabla_V U$\par
            或者是\par
            $\nabla_U: V \mapsto  \nabla_U V$\par
            $\nabla_V: U \mapsto  \nabla_V U$\par
            由于Levi-Civita联络还可以确保在平行移动过程中的无挠(torsion free)，那么还有具有对易性的Lie括号诱导出的（方向导数）算子：\par
            $TM \times TM \to TM$\par
            $(U,V) \mapsto [U,V] = \nabla_U V - \nabla_V U$\par
            定义Riemann曲率张量 $R(U,V)$ 对向量场 $W$ 的作用为：\par
            $R(U,V)W = (\nabla_U\nabla_V - \nabla_V\nabla_U - \nabla_{[U,V]})W$\par
            微分几何是涉及到大量指标计算的领域，这里先不展开，仅仅对上式简单说明．所谓联络的\textbf{曲率(curvature)}\index{curvature}，它实际上给出了协变导数不对易的程度．我们只看上式中的RHS的算子部分：\par
            $\nabla_U\nabla_V - \nabla_V\nabla_U - \nabla_{[U,V]} = (\nabla_U\nabla_V - \nabla_V\nabla_U) - \nabla_{[U,V]}$\par
            前两项相当于是某种形式的对易子，可以理解为：\par
            $[\nabla_U, \nabla_V = \nabla_U\nabla_V - \nabla_V\nabla_U$\par
            另外一项 $ - \nabla_{[U,V]}$ 则是对这个对易子的修正．\par

            实际上上面定义的Riemann曲率张量是一个 $(1,3)$ 型张量，（请回忆流形上张量场的定义 [MP23：张量积、张量、张量丛](https://zhuanlan.zhihu.com/p/45703308)）三个逆变元就是这里的 $U,V,W \in TM$ ，此外还有一个协变元 $\mu \in T^*M$ ，张量对 $(1,3)$ 型输入的作用为：\par
            $\mu \big( R(U,V)W \big)$\par
            我们用一个坐标系来具体讨论，令切向量的基为 $\{\partial_k\}$ ，注基本身也是切向量，于是\par
            $R(\partial_i,\partial_j)\partial_k = (\nabla_i\nabla_j - \nabla_j\nabla_i)\partial_k$\par
            $= \nabla_i\nabla_j\partial_k - \nabla_j\nabla_i\partial_k$\par
            根据前面的讨论，最终可以用Christoffel记号把这个 $(1,3)$ 型张量展开为多重线性形：\par
            $R_{ijk}^l = \partial_i\Gamma_{jk}^l - \partial_j\Gamma_{ik}^l + \Gamma_{jk}^m\Gamma_{im}^l - \Gamma_{ik}^m\Gamma_{jm}^l$\par
            细节——密密麻麻的指标运算，就不给了．\par
            Riemann曲率张量是个 $(1,3)$ 型张量，它的缩并为：\par
            $R_{ij} = R_{ikj}^k$\par
            这个 $(0,2)$ 型张量称为\textbf{Ricci张量(Ricci tensor)}\index{Ricci tensor}，进一步有：\par
            $R = R_i^i$\par
            称为Ricci标量或者\textbf{标量曲率(scalar curvature)}\index{scalar curvature}．这样我们可以构建重要的\textbf{Einstein张量(Einstein tensor)}\index{Einstein tensor}：\par
            $G_{ij} = R_{ij} - \frac{1}{2}Rg_i^j$\par
            它是对空间弯曲情况的描述．\par

    \section{Hodge星算子}
            令$V\in\textbf{Vct}_\mathbb R$为$n$-维实线性空间．$\wedge^pV$上有基$\{\sigma^I=\sigma^{i_1} \wedge \cdots \wedge \sigma^{i_p}\}$．\textbf{体积形式(volume form)}\index{volume form}为：\par
            \[
                \omega = \sigma^1\wedge\cdots\wedge\sigma^n \in \wedge^n V \\
            \]
            $\alpha \in \wedge^pV$如果可以写为\par
            \[
                \alpha = \alpha^1\wedge\cdots\wedge\alpha^p,\quad \alpha^i\in V \\
            \]
            的形式，则称$\alpha$为\textbf{可分解的(decomposable)}\index{decomposable}．\par
            给定$\lambda \in \wedge^p V$，对于任意$\theta \in \wedge^{n-p}V$，$\lambda \wedge \theta \in \wedge^n V$．由于$n$-维线性空间的$n$-形式都是体积形式$\omega \in \wedge^n V$的倍数，可以定义函数：\par
            \begin{align}
                f_\lambda: \wedge^{n-p}V &\to \mathbb R \nonumber \\
                \theta &\mapsto f_\lambda(\theta) \nonumber \\
            \end{align}
            使得\par
            \[
                f_\lambda(\theta)\omega = \lambda \wedge \theta \\
            \]
            如此定义的函数$f_\lambda$是一个线性函数，根据Riesz定理，在内积$g$下它唯一对应了$\varphi \in \wedge^{n-p}V$使得\par
            \[
                f_\lambda(\theta) = g(\theta,\varphi) \\
            \]
            于是定义\textbf{Hodge星算子(Hodge star opertor)}\index{Hodge star opertor}为\par
            \begin{align}
                *: \wedge^pV &\to \wedge^{n-p}V \nonumber \\
                \lambda &\mapsto *\lambda=(-1)^s\varphi \nonumber \\
            \end{align}
            即\par
            \[
                \lambda \wedge \theta = (-1)^sg(\theta,*\lambda)\omega \\
            \]

现在讨论三维空间向量\textbf{叉乘} $\mathbf{u} \times \mathbf{v}$的本质：

$$
\mathbf{u} \times \mathbf{v} =
\begin{pmatrix}
u_2 v_3 - u_3 v_2 \\
u_3 v_1 - u_1 v_3 \\
u_1 v_2 - u_2 v_1
\end{pmatrix} =
\begin{vmatrix}
\mathbf{i} & \mathbf{j} & \mathbf{k} \\
u_1 & u_2 & u_3 \\
v_1 & v_2 & v_3
\end{vmatrix} 
$$

其中 $\mathbf{i}, \mathbf{j}, \mathbf{k}$ 是标准正交基向量．

叉乘是二元线性运算，
暗示了外微分的性质：
\begin{itemize}
    \item $u \times v = -v \times u$
    \item $u \times v = 0$
\end{itemize}

而叉乘只能在三维空间定义．
实际上$u$和$v$并非是向量，
而是$1$-形式，
叉乘其实是外积$\wedge$，
因而$u \wedge v$是三维空间中的$2$-形式．
为何$u \times v$具有向量的形式，
本质是因为

\begin{framed}
\noindent
$n$-维定向Riemann流形$(X,g)$上，
存在映射：

$$
\begin{aligned}
  *: \Omega^k(X) &\to \Omega^{n-k}(X) \\
  \omega &\mapsto *\omega
\end{aligned}
$$

使得$\forall \alpha \in \Omega^k(X)$满足

$$
\alpha \wedge (*\beta) = \langle \alpha, \beta \rangle \, \text{vol}_g, \qquad 
$$

其中：
\begin{itemize}
    \item $\langle \alpha, \beta \rangle$ 是度量 $g$ 诱导的 $k$-形式上的内积，
    \item $\text{vol}_g = \sqrt{\det g}\, dx^1 \wedge \cdots \wedge dx^n$ 是黎曼体积形式．
\end{itemize}
\end{framed}

\subsection{三维空间中的对应关系}

设 $\mathbb{R}^3$ 具有标准欧氏度量与定向．通过 Hodge 星算子 $*$，向量场与微分形式之间有以下对应：
\[
\begin{aligned}
\text{向量} \quad &\longleftrightarrow \quad 1\text{-形式} \\
\mathbf{v} = (v_1,v_2,v_3) \quad &\longleftrightarrow \quad v^\flat = v_1 dx + v_2 dy + v_3 dz,
\end{aligned}
\]
其中 $\flat$ 表示由度量诱导的“降指标”映射（在标准欧氏度量下就是取相同的分量系数）．同样，$2$-形式可通过 $*$ 映射为 $1$-形式：
\[
*: \Omega^2(\mathbb{R}^3) \to \Omega^1(\mathbb{R}^3).
\]

\subsection{与向量分析的关系}

\begin{itemize}
    \item 叉乘的 Hodge 定义将向量积统一到外代数的框架中，无需依赖右手法则的显式公式．
    \item 旋度 $\nabla\times\mathbf{F}$ 可视为 $*(d(\mathbf{F}^\flat))$：
    \[
    \nabla\times\mathbf{F} = \left( * d(\mathbf{F}^\flat) \right)^\sharp.
    \]
    \item 散度 $\nabla\cdot\mathbf{F}$ 对应 $* d( * (\mathbf{F}^\flat) )$．
\end{itemize}

\subsection{总结}

Hodge 星算子与外积共同给出了叉乘的几何定义：
\[
\boxed{\mathbf{u}\times\mathbf{v} = \left( *(\mathbf{u}^\flat \wedge \mathbf{v}^\flat) \right)^\sharp }.
\]
该表述揭示了叉乘实质上是 $2$-形式到 $1$-形式的 Hodge 对偶，从而将三维向量的叉乘整合到外微分形式的一般理论中．



\subsection{叉乘的 Hodge 定义}

对向量 $\mathbf{u}, \mathbf{v} \in \mathbb{R}^3$，定义其叉乘 $\mathbf{u}\times\mathbf{v}$ 为唯一满足下式的向量：
\[
(\mathbf{u}\times\mathbf{v})^\flat = *(\mathbf{u}^\flat \wedge \mathbf{v}^\flat).
\]
即：先取对应的 $1$-形式 $\mathbf{u}^\flat$ 和 $\mathbf{v}^\flat$，作外积得 $2$-形式 $\mathbf{u}^\flat \wedge \mathbf{v}^\flat$，再用 Hodge 星算子将其转化为 $1$-形式，最后通过逆映射 $\sharp$ 变回向量．

\subsection{具体计算验证}

设 $\mathbf{u}=u_1dx+u_2dy+u_3dz$，$\mathbf{v}=v_1dx+v_2dy+v_3dz$（为简便记法，省略 $\flat$ 符号）．则
\[
\mathbf{u}\wedge\mathbf{v} = (u_1v_2-u_2v_1)\,dx\wedge dy + (u_2v_3-u_3v_2)\,dy\wedge dz + (u_3v_1-u_1v_3)\,dz\wedge dx.
\]

Hodge 星算子作用：
\[
\begin{aligned}
*(dx\wedge dy) &= dz, \\
*(dy\wedge dz) &= dx, \\
*(dz\wedge dx) &= dy.
\end{aligned}
\]

因此
\[
*(\mathbf{u}\wedge\mathbf{v}) = (u_2v_3-u_3v_2)\,dx + (u_3v_1-u_1v_3)\,dy + (u_1v_2-u_2v_1)\,dz,
\]
其分量恰为 $\mathbf{u}\times\mathbf{v}$ 的分量．再通过 $\sharp$ 映射得：
\[
\mathbf{u}\times\mathbf{v} = (u_2v_3-u_3v_2,\; u_3v_1-u_1v_3,\; u_1v_2-u_2v_1).
\]



\section{Stokes定理}

            现在，我们要把拓扑和微分形式放在一起，学习\textbf{流形上的Stokes定理}\index{流形上的Stokes定理}，也为进一步学习\textbf{上同调}\index{上同调}作准备．\par

            微积分基本定理是Newton-Leibniz定理：\par
            $\int_{x_1}^{x_2} f^\prime(x) dx = f(x_2) - f(x_1)$\par
            从初学微积分的角度看：N-L定理联系了导数和定积分的关系．\par
            如今我们学了很多的拓扑和微分几何，从现在的角度看：N-L定理联系了局部变化 $f^\prime dx$ 和全局变化 $f(x_2) - f(x_1)$ 的关系．\par
            N-L定理可以有非常直观的理解．假设有一根水管，其位置坐标为 $x\in[x_1,x_2$] ，水管处处漏水而有稳定的流速，于是水管在不同位置坐标上的流速构成函数 $f(x)$ ．\par
            在水管的两端可以测得流速分别为 $f(x_1)$ 和 $f(x_2)$ ，那么在整个区间 $[x_1,x_2$] 上的总漏水速率为：\par
            $f(x_2) - f(x_1)$\par
            这种测量方法只考虑边界两个端点上的取值，是全局的．另一种计算总漏水速率的方式是在每一个位置 $x$ 上测漏水速率，也就是流速的变化量 $ f^\prime(x)$ ，然后在区间上进行积分，得到\par
            $\int_{x_1}^{x_2} f^\prime(x) dx$\par
            这种测量方法是要对整个集合进行逐点操作的，是局部的．\par
            N-L定理揭示了一个事实：\par
            > \textbf{变化在区间上的积分等于边界上呈现的总体变化}\index{变化在区间上的积分等于边界上呈现的总体变化}\par

            电磁学的一个基本定律是\textbf{Gauss静电场定律}\index{Gauss静电场定律}：令 $F=F^xi+F^yj+F^zk$ ，那么在空间闭区域 $\Omega$ 上进行积分，对区域进行如下体积分等于对闭曲面边界 $\partial\Omega$ 进行相应的面积分：\par
            $\iiint_\Omega (\nabla \cdot F)dxdydz = \iint_{\partial \Omega} F^xdydz + F^ydzdx + F^zdxdy$\par
            这个公式的物理意义非常明确．如果 $F$ 为电场，那么左边的被积函数为散度 $\nabla \cdot F$ ，按照面积元 $dxdydz$ 对整个区域积分，得到区域 $\Sigma$ 中总的电荷源．右边的被积表达式表示穿过边界 $\partial\Sigma$的电场，它和区域内的电源是相等的．\par
            更直观地说，对于区域 $\Sigma$ 内部的一个点电荷，其电场是从这一点向外\textbf{发散}\index{发散}的，这个发散的程度就叫做\textbf{散度}\index{散度}；区域 $\Sigma$ 内所有点上都可以有一个散度，它们线性地叠加起来，成为在整个区域中电场（切确地说是 $3$ -微分形式）的变化量．和N-L定律中的区间积分类似，这种计算方法是逐点局部处理的，在中学和普通物理中，我们经常听到所谓的微元法，取一个体积元计算电场散度，这个体积元就是$3$ -微分形式 $dxdydz = dx \wedge dy \wedge dz$ ．\par
            对于区域 $\Sigma$ 内部的一个点电荷，无论用什么曲面去包围它，只要曲面把它包围起来，曲面是闭的，曲面同胚与球面，那么这个点电荷所\textbf{散发}\index{散发}的电场就都会穿过曲面，产生相同的电\textbf{通量}\index{通量}．我们注意到对于曲面而言这是一个拓扑性质，这就是为什么我们前面要学那么多的拓扑打基础．点电荷\textbf{散发}\index{散发}的电场，以\textbf{散度}\index{散度}的方式对应着体积元，而这个点电荷贡献出来的电场，会反映到曲面的电通量上，而\textbf{通量}\index{通量}二字反映的是通过边界曲面的量．和N-L定律中的边界求差相似，通量在边界曲面上的积分也是全局求差的过程，求得的是边界内外之间的差．\par
            Stokes定理将推广这个事实．下面先从拓扑开始，建立最简单的标准单形上的Stokes定理，然后进一步将它推广到微分流形．\par

            前面是以单形为基础建立同调群的，这是\textbf{单纯同调(simplicial homology)}\index{simplicial homology}的处理范围．这种结构极为规整，易于建立积分的概念．\par
            $\mathbb{R}^n$ 中的$n$ \textbf{-标准单形(standard simplex)}\index{standard simplex} $\sigma$ 是坐标原点 $p_0=(0, \dots, 0)$ 和各坐标轴上的单位向量点 $p_k$ 顺序成的，其中\par
            $p_k=(x_0, \dots, x_i, \dots, x_n), \ x_i = \delta_i^k$\par
            例如三维空间的标准 $3$ -单形 $[p_0p_1p_2p_3$] 的点列为：\par
            $p_0=(0,0,0)$\par
            $p_1=(1,0,0)$\par
            $p_2=(0,1,0)$\par
            $p_3=(0,0,1)$\par
            $\mathbb{R}^n$ 中的 $n$ -微分形式是\textbf{体积元(volume form)}\index{volume form}，可以表示为：\par
            $\omega = W(x) dx^1 \wedge dx^2 \wedge \dots \wedge dx^n$\par
            在标准单形 $\sigma$ 这样规整的结构上，借助重积分定义\par
            $\int_\sigma \omega = \int \dots \int_\sigma W(x)dx^1dx^2\dots dx^n$\par
            我们第一次严格地建立起了微分形式的积分．\par

            $\mathbb{R}^n$ 上的 $(n-1)$ 形式\par
            $\psi = \Psi(x)dx^1 \wedge \dots \wedge dx^{n-1}$\par
            在 $n$ 标准单形 $\sigma$ 上有\textbf{Stokes定理}\index{Stokes定理}：\par
            $\int_{\partial\sigma} \psi = \int_\sigma d\psi$\par
            详细的证明篇幅较大，请参考\cite{Nakahara2003}．大约的证明思想如下：\par
            对$(n-1)$ 形式\par
            $\psi = \Psi(x)dx^1 \wedge \dots \wedge dx^{n-1}$\par
            进行外微分的结果是一个 $n$ 形式\par
            $d\psi = \text{psi}(x)dx^1 \wedge \dots \wedge dx^n$\par
            它是一个体积元，于是根据前述定义其微分形式的积分是一个 $n$ 重积分：\par
            $\int_\sigma d\psi = \int \dots \int_\sigma \text{psi}(x)dx^1\dots dx^n$\par
            再考虑边界 $\partial\sigma$ 上的积分． $n$ 单形 $\sigma$ 的边界是 $(n-1)$ 单形，于是\par
            $\int_{\partial\sigma} \psi = \int \dots \int_{\partial\sigma} \Psi(x)dx^1\dots dx^{n-1}$\par
            也成为$(n-1)$ 单形 $\partial \sigma$ 上的$(n-1)$ 形式 $\psi$ 的积分，它是 $(n-1)$ 重积分．\par
            经过两边的计算，两个重积分是相等的，于是得证．\par
            在标准单形中代入Newton-Leibniz公式或者Gauss电场公式，都可以验证标准单形的Stokes定理．下篇将把这个定理推广到微分流形上．\par

            为了把Stokes定理拓展到微分流形，我们需要先微分形式的积分从标准单形推广到微分流形上．我们现在要建立几何上更加灵活的\textbf{奇异同调(singular homology)}\index{singular homology}，它保持了单纯同调的拓扑特性，但不再被限制于单纯同调中的三角剖分．奇异同调建立的基础是由标准单形\textbf{光滑同胚(diffeomorphism)}\index{diffeomorphism}映射而成的奇异单形：\par
            $f:\sigma \to s$\par
            由于是光滑同胚，显然保持了拓扑性质．（注：实际上只需要光滑的条件就可以解决后继问题，我们这里加强为同胚的条件是为了保持单形的拓扑性质，便于直观理解）光滑映射 $f$ 的灵活性确保了从标准单形 $\sigma$ 出发可以得到任意所需的（具有相同拓扑性质的）光滑流形 $s$ ，相对于 $n$ -标准单形 $\sigma$ ，称 $s$ 为 $n$ -\textbf{奇异单形(singular simplex)}\index{singular simplex}．\par
            以奇异单形为基础，通过类似于 [MP22：同调群](https://zhuanlan.zhihu.com/p/44776557) 的构造方式，可以得到\textbf{奇异同调群(singular homological group)}\index{singular homological group}：\par
            $ H_k(C) = Z_k/B_k = \text{Ker}(\partial_k)/ \text{Im}(\partial_{k+1})$\par
            由于我们加强了光滑同胚的条件，这样得到的单纯同调群和奇异同调群并无拓扑性质的区别，我们在符号和阐述上也不区分二者．建立奇异同调群的意义在于，我们用单纯同调的方法就可以得到任意微分流形的拓扑性质，于是我们在前面大量阐述的拓扑知识，现在可以不加区分地用于微分流形：\par

            为了把微分形式的积分从标准单形推广到微分流形，我们首先要建立奇异单形上的积分．光滑同胚映射\par
            $f:\sigma \to s$\par
            的作用相当于把积木一样的标准单形捏成了任意形状，倘若我们直接定义奇异单形上的积分：\par
            $\int_s \omega = \int_\sigma \omega = \int \dots \int_\sigma w(x)dx^1dx^2\dots dx^n$\par
            是一定有问题的，它没有反映出映射 $f$ 的效果，若要在微分形式上反映出映射的效果，需要用到拉回映射．微分流形间的光滑映射\par
            $f:M \to N$\par
            $p \mapsto f(p)$\par
            可以诱导出两个点的余切空间之间的\textbf{拉回映射(pullback)}\index{pullback}：\par
            $f^*: T_{f(p)}^*N \to T_p^*M$\par
            注意拉回映射 $f^*$ 的方向与 $f$ 相反，以此得名．也就是当$f$把 $p$ 映射为 $f(p)$ 时，拉回映射 $f^*$ 也同时把 $f(p)$ 上的余切向量映射回 $p$ 上的余切向量．由于切空间和余切空间对偶， $\forall v \in T_pM, \ \omega \in T_{f(p)}^*N$ 规定以上两个映射的条件是：\par
            $\langle f^*\omega, v \rangle_M = \langle \omega, f_*v \rangle_N$\par
            从而使拉回映射具有以下性质：\par
            $d(f^*\omega) = f^*(d\omega)$\par
            现在，针对标准单形到奇异单形的映射\par
            $f:\sigma \to s$\par
            其拉回为：\par
            $f^*: T_{f(p)}^*s \to T_p^*\sigma$\par
            我们定义\textbf{奇异单形上微分形式的积分}\index{奇异单形上微分形式的积分}为：\par
            $\int_s \omega = \int_\sigma f^*\omega$\par

            经过前面建立的奇异同调群和奇异单形上微分形式的积分，通过类似 [MP22：同调群](https://zhuanlan.zhihu.com/p/44776557) 中的形式和，不难定义奇异链 $c = \sum_k z_ks_k \in C_r(M)$ 上的积分\par
            $\int_c \omega = \sum_k z_k \int_{s_k} \omega$\par
            令微分形式 $w \in \Omega^{n-1}(M)$ ，奇异链 $c \in C_r(M)$ ，则有\textbf{Stokes定理}\index{Stokes定理}：\par
            $\int_{\partial c} \omega = \int_c d\omega$\par
            \textbf{证明}\index{证明}：\par
            考虑 $n$ -奇异单形 $s$ ．根据奇异单形上微分形式的积分定义，以及拉回映射的性质：\par
            $\int_c d\omega = \int_\sigma d(f^*\omega) = \int_\sigma f^*(d\omega)$\par
            以及\par
            $\int_{\partial c} \omega = \int_{\partial\sigma} f^*\omega$\par
            前面讲到， $\mathbb{R}^n$ 上的 $(n-1)$ 形式 $\psi$ ，在 $n$ 标准单形 $\sigma$ 上有Stokes定理：\par
            $\int_{\partial\sigma} \psi = \int_\sigma d\psi$\par
            于是\par
            $\int_{\partial\sigma} f^*\omega = \int_\sigma d(f^*\omega)$\par
            证明了\textbf{奇异链上的Stokes定理}\index{奇异链上的Stokes定理}：\par
            $\int_{\partial c} \omega = \int_c d\omega$\par
            奇异链上的Stokes定理的特例是\textbf{（定向）流形上的Stokes定理}\index{（定向）流形上的Stokes定理}：令 $D$ 是定向微分流形 $M$ 的开子流形． $M$ 诱导出 $D$ 和 $\partial D$ 的定向．令 $\omega$ 是 $M$ 上的 $(n-1)$ 次光滑形式：\par
            $\int_D d\omega = \int_{\partial D} \omega$\par
            Newton-Leibniz、Green、Gauss、Stokes定理，在微分几何中都可以归纳到这个形式中．\par
            它概括了刚才在N-L定理和Gauss静电场定律中揭示的事实：\par
            > \textbf{变化在区间上的积分等于边界上呈现的总体变化}\index{变化在区间上的积分等于边界上呈现的总体变化}\par
            引入类似内积的符号：\par
            $\int_D d\omega = \langle D, d\omega \rangle$\par
            $\int_{\partial D} \omega = \langle \partial D, \omega \rangle$\par
            于是\par
            $\langle D, d\omega \rangle = \langle \partial D, \omega \rangle$\par
            这里暗示了整体算子 $\partial$ 和局部算子 $d$ 之间的对偶关系．我们现在已经完成了对偶性的一面，即在\textbf{同调(homology)}\index{homology}和边界算子 $\partial$ 方面的研究．下一步是对偶性的另一面，\textbf{上同调(cohomology)}\index{cohomology}和外微分算子 $d$ 的研究．\par


\section{外微分 Stokes 定理的几种常见形式}

\subsection{统一形式（外微分）}

设 $M$ 为 $n$ 维带边光滑流形，$\omega$ 为 $(n-1)$-形式（紧支集），则
\[
\int_{\partial M} \omega = \int_M d\omega.
\]
以下各种定理均为其特例．

\subsection{微积分基本定理（$n=1$）}

取 $M=[a,b]\subset\mathbb{R}$，$\omega = f$（$0$-形式），则 $d\omega = f'(x)\,dx$，$\partial M = \{b\}-\{a\}$（定向相反）：
\[
\int_a^b f'(x)\,dx = f(b)-f(a).
\]

\subsection{格林定理（$n=2$，平面区域）}

设 $M\subset\mathbb{R}^2$ 为平面区域，$\omega = P\,dx + Q\,dy$（$1$-形式），则
\[
d\omega = \left(\frac{\partial Q}{\partial x} - \frac{\partial P}{\partial y}\right) dx\wedge dy.
\]
Stokes 定理给出：
\[
\int_{\partial M} P\,dx + Q\,dy = \iint_M \left(\frac{\partial Q}{\partial x} - \frac{\partial P}{\partial y}\right) dx\,dy.
\]

\subsection{经典 Stokes 定理（$n=2$，三维空间中的曲面）}

设 $M\subset\mathbb{R}^3$ 为二维定向曲面，$\mathbf{F}=(F_x,F_y,F_z)$ 为向量场，对应 $1$-形式 $\omega = F_x dx + F_y dy + F_z dz$．则 $d\omega$ 对应于旋度 $\nabla\times\mathbf{F}$ 的 $2$-形式：
\[
\int_{\partial M} \mathbf{F}\cdot d\mathbf{r} = \iint_M (\nabla\times\mathbf{F})\cdot d\mathbf{S}.
\]

\subsection{散度定理（$n=3$，三维区域）}

设 $M\subset\mathbb{R}^3$ 为三维区域，定义 $2$-形式
\[
\omega = F_x\,dy\wedge dz + F_y\,dz\wedge dx + F_z\,dx\wedge dy.
\]
则 $d\omega = (\nabla\cdot\mathbf{F})\,dx\wedge dy\wedge dz$，Stokes 定理给出：
\[
\int_{\partial M} \mathbf{F}\cdot d\mathbf{S} = \iiint_M (\nabla\cdot\mathbf{F})\,dV.
\]

\subsection{总结对照表}

\[
\begin{array}{|c|c|c|c|c|}
\hline
\text{名称} & n & \omega & d\omega & \text{表达式} \\ \hline
\text{微积分基本定理} & 1 & f & f'(x)dx & \int_a^b f'(x)dx = f(b)-f(a) \\ \hline
\text{格林定理} & 2 & Pdx+Qdy & (Q_x-P_y)dx\wedge dy & \int_{\partial M} Pdx+Qdy = \iint (Q_x-P_y)dxdy \\ \hline
\text{经典 Stokes} & 2 & \mathbf{F}\cdot d\mathbf{r} & (\nabla\times\mathbf{F})\cdot d\mathbf{S} & \int_{\partial M}\mathbf{F}\cdot d\mathbf{r} = \iint (\nabla\times\mathbf{F})\cdot d\mathbf{S} \\ \hline
\text{散度定理} & 3 & \mathbf{F}\cdot d\mathbf{S} & (\nabla\cdot\mathbf{F})\,dV & \int_{\partial M}\mathbf{F}\cdot d\mathbf{S} = \iiint (\nabla\cdot\mathbf{F})\,dV \\ \hline
\end{array}
\]

所有定理统一于外微分形式下的 Stokes 公式 $\int_{\partial M}\omega = \int_M d\omega$．

\chapter{de Rham上同调}

\section{外微分算子}

在光滑流形$X$上，
微分形式代数$(\Omega^\bullet(X),\wedge,\text d)$就是这样的外代数，
其中的外微分算子

$$
\text d^i: \Omega^i(X) \to \Omega^{i+1}(X)
$$

提升了外微分的阶数，
特别是最初的$f \mapsto \text d f$。
我们熟知

$$
\text d \circ \text d = 0
$$

这意味着在态射序列中：

$$
\Omega^0(X) \stackrel{\text d^0}{\longrightarrow} 
\Omega^1(X) \stackrel{\text d^1}{\longrightarrow} 
\Omega^2(X) \stackrel{\text d^2}{\longrightarrow} \cdots
$$

所有阶数在$\Omega^{i+1}(X)$上都满足$\text{Im} \ \text d^i \subseteq \text{Ker} \ \text d^{i+1}$。
这样便可以用正合序列的方法来研究外微分形式代数$\Omega^\bullet(X)$。


\chapter{概型}

\section{Zariski拓扑}

前面由子集生成理想的过程是：
$\textbf 2^A \to I_A :: S \mapsto (S)$。
现在限制在素谱$\text{Spec} \ A \subset I_A$上，
$V: \textbf 2^A \to I_A :: S \mapsto V(S)$。


给定交换环的子集$E \subset A \in\textbf{CRing}$，
一个理想$I \subset E$去覆盖子集的意义。

  上述的例子中，最终被素理想所覆盖，
可见素理想的重要性。记
$$
V(E) = \left\{ \mathfrak p \in X \ | E \subset \mathfrak p \right\}
$$
为\textbf{退化集(vanishing set)}\index{vanishing set 退化集}。


退化集$V(E)$由包含子集$E$的素理想构成。
显然$E$生成的理想$(E)$满足
$$
V((E)) = V(E)
$$
并且易于证明：
\begin{enumerate}
\item $V(0) = X$，$V(1)=\varnothing$
\item 对于任意子集族$\{E_i\in 2^A\}_{i\in I}$有$V(\cup_{i\in I} E_i) = \bigcap_{i\in I} V(E_i)$
\item 对于$A$的任意理想$\mathfrak a$和$\mathfrak b$有$V(\mathfrak a \cap \mathfrak b) = V(\mathfrak a \mathfrak b) = V(\mathfrak a) \cup V(\mathfrak b)$
\end{enumerate}
以上条件使得退化集符合闭集的公理，
于是退化集作为闭集赋予了$X=\text{Spec} \ A$拓扑空间的性质，
称为\textbf{Zariski拓扑(Zariski topology)}\index{Zariski topology Zariski拓扑}。

\begin{example}

素数$p$生成素理想$(p) \in \text{Spec} \ \mathbb Z$，
此外还有一个素理想$(0) \in \text{Spec} \ \mathbb Z$称为
\textbf{泛点(generic point)}\index{generic point 泛点}。
如下表示以示区分：
$$
\text{Spec} \ \mathbb Z = (0) \cup \{(p) \ | \ p > 1\}
$$
其中$p$为素数。
\end{example}

\begin{example}
令$E=\{3,4,5,6\}$则$(E)=(60)$，
$V(E)=V(60)$由包含$(60)$的素理想构成，
小于$60$的素数所生成的理想都是$V(E)$中的点。
\end{example}


\section{多项式环}
\begin{framed}
\noindent\textbf{定义（多项式环）：}设 $k$ 为域，由 $n$ 个变量和系数域 $k$ 生成的环 $R = k[x_1, \dots, x_n]$ 称为\textbf{多项式环}。
\end{framed}
作为最具代表性的交换环，它是所有代数几何现象发生的基础舞台。

\section{Zariski拓扑：用理想构造}
不同于经典代数几何研究方程的具体坐标零点，格罗滕迪克将其代数化：将环 $R$ 的\textbf{所有素理想}构成的集合定义为仿射空间，记作 $\operatorname{Spec}(R)$。
对于任意理想 $I \subseteq R$，定义闭集为包含 $I$ 的所有素理想：
$$ V(I) = \{ \mathfrak{p} \in \operatorname{Spec}(R) \mid I \subseteq \mathfrak{p} \} $$
由这些闭集生成的拓扑，即为定义在 $\operatorname{Spec}(R)$ 上的 Zariski 拓扑。

\section{概型}
\begin{framed}
\noindent\textbf{定义（仿射概型与概型）：}将局部赋环空间的几何直觉与素理想谱 $\operatorname{Spec}(R)$ 的代数构造融合，得到具有结构层 $\mathcal{O}_{\operatorname{Spec}(R)}$ 的局部赋环空间，称为\textbf{仿射概型}。
\textbf{概型}（Scheme）则是一个更为一般的局部赋环空间，它的每一点都存在一个开邻域，该邻域作为一个诱导的局部赋环空间同构于某个仿射概型。
\end{framed}
概型是现代代数几何探讨几何本质的最核心语言。

\section{有理函数域}
\begin{framed}
\noindent\textbf{定义（泛点与有理函数域）：}对于一个由整环生成的积分概型，空间中存在一个特殊的“\textbf{泛点}”（Generic Point），它对应于环的零理想 $(0)$。结构层在这个泛点处的茎 $\mathcal{O}_{X, (0)}$ 恰好同构于整环的分式域，即该概型的\textbf{有理函数域} $k(X)$。
\end{framed}
在几何上，这代表了整个空间上所有“除了在某些低维闭子集之外，几乎处处有定义”的函数之全体。
\chapter{辛几何}

\section{近复结构}

复数可以如下表示，
体现实线性映射、酉群$\text{U}(1)\simeq \text{SO}(2)$：

$$
z = x + y \text i 
= |z| \cdot \text e^{\theta \text i} = \begin{bmatrix}
  x & -y \\ y & x 
\end{bmatrix}
= \sqrt{x^2 + y^2} \begin{bmatrix}
  \cos\theta & -\sin\theta \\ \sin\theta & \cos\theta
\end{bmatrix}
$$

若把$\mathbb C$视为$2 \times 2$实矩阵空间$\mathbb R_2^2 \in\textbf{Vct}_\mathbb R$中的一个二维线性子空间，
取基：

$$
I = \begin{bmatrix}
  1 & 0 \\ 0 & 1
\end{bmatrix}, \quad
J = \begin{bmatrix}
  0 & -1 \\ 1 & 0
\end{bmatrix}, \quad
$$

使得$z = x I + y J$，
且有$J^2 = J \circ J = -I$，
相当于$\text i^2 = -1$。

现在延伸到$n$-维复矩阵空间：

$$
\big( \mathbb C_n^n \in\textbf{Vct}_\mathbb C \big)
\simeq
\big( \mathbb R_{2n}^{2n} \in\textbf{Vct}_\mathbb R \big)
$$

一个复矩阵类似表示为：

$$
\begin{bmatrix}
  z_1^1 & \cdots & z_1^n \\
  \vdots & \ddots & \vdots \\
  z_n^1 & \cdots & z_n^n \\
\end{bmatrix}
=
\begin{bmatrix}
  x_1^1 & -y_1^1 & \cdots & x_1^n & -y_1^n \\
  y_1^1 & x_1^1 & \cdots & y_1^n & x_1^n \\
  \vdots & \vdots & \ddots & \vdots & \vdots \\
  x_n^1 & -y_n^1 & \cdots & x_n^n & -y_n^n \\
  y_n^1 & x_n^1 & \cdots & y_n^n & x_n^n \\
\end{bmatrix}
$$

为了构造分块矩阵，取它等价的表示形式：

$$
\begin{aligned}
\begin{bmatrix}
  z_1^1 & \cdots & z_1^n \\
  \vdots & \ddots & \vdots \\
  z_n^1 & \cdots & z_n^n \\
\end{bmatrix}
&=
\left[
\begin{array}{ccc|ccc}
  x_1^1 & \cdots & x_1^n & -y_1^1 & \cdots & -y_1^n \\
  \vdots & \ddots & \vdots & \vdots & \ddots & \vdots \\
  x_n^1 & \cdots & x_n^n & -y_n^1 & \cdots & -y_n^n \\
  \hline
  y_1^1 & \cdots & y_1^n & x_1^1 & \cdots & x_1^n \\
  \vdots & \ddots & \vdots & \vdots & \ddots & \vdots \\
  y_n^1 & \cdots & y_n^n & x_n^1 & \cdots & x_n^n \\
\end{array}
\right]
\\
&=
\left[
\begin{array}{c|c}
  X & -Y \\
  \hline
  Y & X
\end{array}
\right]
= X I_n + Y J_n
\end{aligned}
$$


类似前述有基：

$$
I = 
\left[
\begin{array}{c|c}
  I_n & 0 \\
  \hline
  0 & I_n
\end{array}
\right]
, \quad
J = 
\left[
\begin{array}{c|c}
  0 & -I_n \\
  \hline
  I_n & 0
\end{array}
\right]
$$

$I$是$\mathbb R_{2n}^{2n}$的代数单位元，
$J$则满足：

$$
J^2 = J \circ J = 
\left[
\begin{array}{c|c}
  0 & -I_n \\
  \hline
  I_n & 0
\end{array}
\right]
\circ
\left[
\begin{array}{c|c}
  0 & -I_n \\
  \hline
  I_n & 0
\end{array}
\right]
=
\left[
\begin{array}{c|c}
  -I_n & 0 \\
  \hline
  0 & -I_n
\end{array}
\right]
= - I.
$$

复数单位$\text i$和基矩阵$J$具有相同的性质：

\begin{framed}
  
实向量空间$V$上的线性变换$J \in \text{End} \ V$若满足

$$
J^2 = -\text{id}_V
$$

称$J$是一个\textbf{近复结构}(almost complex structure)。

\end{framed}

在实矩阵表示下，
实转置对应着复共轭转置，
共轭携带的信息不变。
在不混淆的情况下，
把$J$写为转置/共轭的形式：

$$
J = \left[
\begin{array}{c|c}
0 & I_n\\
\hline
-I_n & 0
\end{array}
\right].
$$

它就是辛矩阵。

\section{辛结构}

作用在$k$-线性空间上，
基于左右元的区分，
一对对偶算子$f$作用在左元，
$f^*$作用在右元．
若$\forall x \in X, \beta \in Y^*:\ \langle x,f^* \beta \rangle = \langle f x, \beta \rangle $，
则互为伴随算子．
这里$\langle -,- \rangle$是$X \times X^*$和$Y \times Y^*$上的自然配对。
考虑$X$上的线性自映射，
自伴算子$A = A^\dagger$，
若引入内积则自然变换可以用内积描述为：

$$
% https://tikzcd.yichuanshen.de/#N4Igdg9gJgpgziAXAbVABwnAlgFyxMJZABgBpiBdUkANwEMAbAVxiRAA0QBfU9TXfIRRkATFVqMWbdgD0AVN14gM2PASIjy4+s1aIOivqsEbSY6jqn7ZCnkYHrhpAIzbJekAB1PDOmADmDDAABAC0pACC8mHB3gBOfoGsdsr8akLImq4W7mzevgFBwRHhobGeCYXJSioOGQAsWjm6bAAehqnGjsgAbE0SLfoR7Sm16UR95gNWIACeHWMmKI1Tlh4RwfOjaUvIjdnTHvmJRa2Rm+WVSQs73X0Ha3k+JyEbZ7OXL9ziMFD+8ERQAAzOIQAC2SDIIBwECQmkObAiIGovgARjAGAAFW5CEBBIE4Dog8FIADM1BhSGczRmUVsSmJEMQjWhsMQxBSjKQLMpiGcnNBTIArBS2SIBSTECLWWSJUy+jLEAB2Cl0LAMNhguhoOCUuVIAAcoqQAE5VerNdrdbD9XyobyFTg1Rr9FqdXqGYKqfa2Wboc7Le6bZ7Jc5qYqVf6La6rR7gV6+eHeUaoy6QG7rUSE84fVTwwwsGAPFAIExUUFkSAABYwOhQNiQIuVp3R8AEZIULhAA
\begin{tikzcd}
X \arrow[rr, "A"]                             &  & X                                          &  & x \arrow[rr, maps to]                                                                         &  & Ax                                                            \\
{\langle -,A^* - \rangle} \arrow[u] \arrow[d] &  & {\langle A-,- \rangle} \arrow[u] \arrow[d] &  & {\langle x,A y \rangle} \arrow[u, maps to] \arrow[d, maps to] \arrow[rr, no head, Rightarrow] &  & {\langle A x,y \rangle} \arrow[u, maps to] \arrow[d, maps to] \\
X^*                                           &  & X^* \arrow[ll, "A^*"]                      &  & A y                                                                                           &  & y \arrow[ll, maps to]                                        
\end{tikzcd}
$$

复线性空间的内积是Hermite内积，
它的实部是实内积，
虚部是辛内积。

线性空间$V$上的双线性形式称为斜对称形式，
若满足：

$$
\omega (u,v) + \omega(v,u) = 0
$$

$2n$-维实线性空间$V$配置非退化斜对称形式$\omega$，
构成辛空间$(V,\omega)$。

若$J$是一个满足
$J^2 = -I$ 的线性变换，则称其为一个\textbf{复结构}。在这种情况下，
可由$J$定义一个由

$$
\omega_J(u,v) := \langle Ju,v\rangle
$$

给出的辛形式。特别地，若$V$配备一个内积$\langle\cdot,\cdot\rangle$，
则一个矩阵$A\in \mathrm{GL}(2n,\mathbb R)$称为\textbf{辛矩阵}，若它满足

$$
A^T J A = J,
$$

在\cite{ChrissGinzburg1997}中如此定义： 

\begin{framed}
    
A symplectic structure on \( X \) is a non-degenerate regular (i.e., \( C^\infty \), resp. holomorphic, algebraic) 2-form \( \omega \) such that \( d\omega = 0 \).

\end{framed}

Let \( X = \mathbb{C}^{2n} \) with coordinates \( q_1, \ldots, q_n, p_1, \ldots, p_n \). Then

\[
\omega = dp_1 \wedge dq_1 + \cdots + dp_n \wedge dq_n
\]

is a symplectic structure.
\chapter{力学}

位形空间$X$通常是一个$n$-维光滑流形，
其上的函数集构成光滑函数代数$C(X)$。
切丛$TX$构成物理上的\textbf{速度相空间}(velocity phase space)，
在局部坐标系表示的广义坐标$x^i$下，
有广义速度$\dot q^i = \text d q^i / \text d t$，
二者构成$TX$的局部坐标。
考虑函数$L(q^i,\dot q^i,t)$，
它在时间上的积分称为\textbf{作用量}(action)：

$$
S = \int_0^T L \text d t
$$

\textbf{Hamilton变分原理}(variational principle of Hamilton)要求作用量最小，
这样的极值问题涉及到对泛函求变分使得：

$$
\delta S = \delta \int_0^T L \text d t = 0
$$

需要以独立的方式去扰动$C(X)$的点。
函数求导扰动的是用实数无穷小量$\varepsilon$扰动$x$：

$$
\frac{\text d f}{\text d x} \bigg|_{x = x_0}
= \lim_{\varepsilon \to 0}
\frac{f(x_0 + \varepsilon) - f(x_0)}{\varepsilon}
$$

而泛函更复杂，
如果上述极限的分母仍是实数无穷小量$\varepsilon$，
那么$\varepsilon$扰动$f$需要引入一个函数$g$来产生
$f(x) + \varepsilon g(x)$。
普通的$g$会同时扰动$f$在许多$x$处的函数值，
无法独立出点$x=x_0$对函数变换的贡献。
为了在$x = x_0$点附近扰动，
狄拉克$\delta$函数就很适合：

$$
\delta(x) = \begin{cases}
+\infty, & x=0,\\
0, & x\ne 0,
\end{cases}
\qquad \text{s.t.} \quad
\int_{-\infty}^{\infty}\delta(x)\,dx = 1.
$$

于是$f(x) + \varepsilon \delta$只在$x = x_0$处扰动了$f$。
应用狄拉克$\delta$函数构造\textbf{泛函导数}(functional derivative)：

$$
\frac{\delta F}{\delta f}
= \lim_{\varepsilon\to 0}
\frac{F[f(x)+\varepsilon\delta(x^\prime - x)]-F[f(x)]}{\varepsilon}
$$

现在需要能让作用量最小的Lagrangian $L$，
满足：

$$
\frac{\delta L}{\delta x} = 0
$$

总机械能$E = T + V$守恒的情况下，
动能和势能此消彼长。
考虑到量子力学只能测量平均值，
在时间区间$[0,\tau]$上通过积分得到

$$
\overline T = \frac{1}{\tau}\int_0^\tau \frac{m \dot x^2 \text d t}{2}, \quad 
\overline V = \frac{1}{\tau}\int_0^\tau V[x(t)] \text d t
$$

求二者的变分\cite{Lancaster2014}：

$$
\frac{\delta \overline T[x]}{\delta x(t)} = - \frac{m \ddot x}{\tau}, \quad
\frac{\delta \overline V[x]}{\delta x(t)} = \frac{V^\prime[x(t)]}{\tau}
$$

于是

$$
\frac{\delta (\overline T[x] - \overline V[x])}{\delta x(t)} = 0
$$

这样得到了我们需要的Lagrangian:

$$
L = T - V
$$

满足\textbf{Hamilton最小作用量原理}(Hamilton's principle of least action)。
视坐标和速度为独立变量:

$$
\frac{\delta S}{\delta x(t)} 
= \frac{\delta L}{\delta x(t)} - \frac{\text d}{\text d t} \frac{\delta L}{\delta \dot x(t)} 
= 0
$$

得到Euler-Lagrange方程。


\cite{MarsdenRatiu1994}




\printbibliography

\end{document}